\nonstopmode{}
\documentclass[a4paper]{book}
\usepackage[times,inconsolata,hyper]{Rd}
\usepackage{makeidx}
\makeatletter\@ifl@t@r\fmtversion{2018/04/01}{}{\usepackage[utf8]{inputenc}}\makeatother
% \usepackage{graphicx} % @USE GRAPHICX@
\makeindex{}
\begin{document}
\chapter*{}
\begin{center}
{\textbf{\huge Package `oystermapR'}}
\par\bigskip{\large \today}
\end{center}
\ifthenelse{\boolean{Rd@use@hyper}}{\hypersetup{pdftitle = {oystermapR: Predict and Map Oyster Growth Suitability from Environmental Data}}}{}
\ifthenelse{\boolean{Rd@use@hyper}}{\hypersetup{pdfauthor = {T Tucker}}}{}
\begin{description}
\raggedright{}
\item[Title]\AsIs{Predict and Map Oyster Growth Suitability from Environmental
Data}
\item[Version]\AsIs{0.4.0}
\item[Description]\AsIs{Predicts spatial suitability for oyster growth from environmental
survey data using Analytic Hierarchy Process (AHP) weighted scoring.
Users supply sensor data (ADCP, CTD, bathymetric sonar, sidescan),
specify a target species, and receive per-location suitability scores,
a GeoTIFF heatmap for QGIS, contour lines, and a formatted PDF or HTML
report. Supported species: Ostrea edulis (European Flat Oyster), Magallana
gigas (Pacific Oyster), Crassostrea angulata (Portuguese Oyster),
Ostrea stentina (Denticulate Flat Oyster), and Ostrea lurida (Olympia
Oyster). Includes season-aware scoring, tidal height correction, Bayesian
tolerance parameter updating from field observations, spatial block
cross-validation (Roberts et al. 2017), permutation variable importance,
wave exposure and sediment stability modules, HAB risk and anthropogenic
disturbance scoring with optional live ICES data integration, hybrid
larval dispersal connectivity scoring (union-find Gaussian kernel plus
optional OpenDrift/FVCOM connectivity matrix), and batch multi-species
comparison.}
\item[License]\AsIs{MIT + file LICENSE}
\item[Encoding]\AsIs{UTF-8}
\item[Roxygen]\AsIs{list(markdown = TRUE)}
\item[RoxygenNote]\AsIs{7.3.3}
\item[Imports]\AsIs{dplyr (>= 1.0.0), terra (>= 1.7.18), sf (>= 1.0.0), rlang (>=
1.0.0), cli (>= 3.0.0), stats, utils}
\item[Suggests]\AsIs{ggplot2 (>= 3.4.0), testthat (>= 3.0.0), knitr, rmarkdown (>=
2.20), tinytex, httr (>= 1.4.0), jsonlite (>= 1.8.0), gstat (>=
2.1.0), sp (>= 1.6.0), lubridate (>= 1.9.0), suncalc (>= 0.5.0)}
\item[Depends]\AsIs{R (>= 4.1.0)}
\item[NeedsCompilation]\AsIs{no}
\item[Author]\AsIs{T Tucker [aut, cre]}
\item[Maintainer]\AsIs{T Tucker }\email{tristantucker48@gmail.com}\AsIs{}
\end{description}
\Rdcontents{Contents}
\HeaderA{oystermapR-package}{oystermapR: Predict and Map Oyster Growth Suitability}{oystermapR.Rdash.package}
\aliasA{oystermapR}{oystermapR-package}{oystermapR}
\keyword{internal}{oystermapR-package}
%
\begin{Description}
\code{oystermapR} provides a rule-based, AHP-weighted suitability scoring
pipeline for predicting where conditions are most likely to support oyster
growth. Users supply a tabular dataset of environmental and physical
measurements (collected in the field or from remote sensing), specify a
target species, and receive:
\begin{itemize}

\item{} Per-location suitability scores [0, 1]
\item{} Categorical suitability classes (High / Moderate / Low / Very Low / Excluded)
\item{} Per-variable component scores for diagnosis and reporting
\item{} A GeoTIFF raster file for heatmap visualisation in QGIS
\item{} A QGIS colour-ramp style file (.qml)
\item{} Contour lines for overlay mapping
\item{} A formatted PDF or HTML report

\end{itemize}

\end{Description}
%
\begin{Section}{Supported species}

\Tabular{llll}{
Key & Latin name & Common name & Region \\{}
\code{ostrea\_edulis} & \emph{Ostrea edulis} & European flat oyster & NE Atlantic, North Sea, Mediterranean \\{}
\code{magallana\_gigas} & \emph{Magallana gigas} & Pacific oyster & Cosmopolitan (introduced) \\{}
\code{crassostrea\_angulata} & \emph{Crassostrea angulata} & Portuguese oyster & Iberian Peninsula, W Atlantic \\{}
\code{ostrea\_stentina} & \emph{Ostrea stentina} & Denticulate flat oyster & Mediterranean, Mar Menor \\{}
\code{ostrea\_lurida} & \emph{Ostrea lurida} & Olympia oyster & NE Pacific: Alaska to Baja California \\{}
}


Use \code{\LinkA{list\_species()}{list.Rul.species}} to see all species with their data quality ratings.
\end{Section}
%
\begin{Section}{Typical workflow}


\begin{alltt}library(oystermapR)

# Load your survey data
df <- read.csv("my_survey.csv")

# Run the prediction
result <- predict_oyster(
  data           = df,
  species        = "ostrea_edulis",
  output_geotiff = "oyster_suitability.tif",
  verbose        = TRUE
)

# Inspect high-suitability sites
subset(result, suitability_class == "High")

# Export matching QGIS colour style
export_qml_style("oyster_suitability.tif")
\end{alltt}

\end{Section}
%
\begin{Section}{Sensor compatibility}

The package is designed around data collectable with:
\begin{itemize}

\item{} \strong{Nortek Signature 500 ADCP} \bsl{}u2014 current velocity, shear stress, turbidity
proxy (read with \code{\LinkA{read\_nortek\_adcp()}{read.Rul.nortek.Rul.adcp}})
\item{} \strong{Ping 3DSS iDX450 PRO Bathymetric Sonar} \bsl{}u2014 depth, slope, roughness
(read with \code{\LinkA{read\_sonar\_tif()}{read.Rul.sonar.Rul.tif}})
\item{} \strong{Lowrance + BioBase} \bsl{}u2014 depth, bottom hardness, vegetation, sidescan
\item{} \strong{Salinity/temperature probes} \bsl{}u2014 temperature, salinity, dissolved oxygen

\end{itemize}

\end{Section}
%
\begin{Section}{Validation and model diagnostics}

Validation functions for research-grade use:
\begin{itemize}

\item{} \code{\LinkA{validate\_against\_records()}{validate.Rul.against.Rul.records}}: AUC, TSS, Brier score, F1, sensitivity,
specificity vs known presence/absence records.
\item{} \code{\LinkA{spatial\_block\_cv()}{spatial.Rul.block.Rul.cv}}: Spatial block cross-validation (Roberts et al. 2017)
\bsl{}u2014 avoids inflated AUC from spatially autocorrelated data.
\item{} \code{\LinkA{permutation\_importance()}{permutation.Rul.importance}}: Variable importance by AUC drop after
permuting each scored variable.
\item{} \code{\LinkA{sensitivity\_analysis()}{sensitivity.Rul.analysis}}: Partial dependence curve \bsl{}u2014 suitability vs a
single variable while others are held at their median.

\end{itemize}

\end{Section}
%
\begin{Section}{Bayesian tolerance updating}

Tolerance parameters (optimal ranges) can be updated from field observations:
\begin{itemize}

\item{} \code{\LinkA{update\_species\_tolerances()}{update.Rul.species.Rul.tolerances}}: MAP + Laplace approximation (fast) or
RWMH MCMC (full posterior). Sequential \bsl{}u2014 posterior becomes next prior.
\item{} \code{\LinkA{get\_tolerance\_posteriors()}{get.Rul.tolerance.Rul.posteriors}}: Retrieve updated parameters with 95\% CIs.
\item{} \code{\LinkA{save\_tolerance\_update()}{save.Rul.tolerance.Rul.update}} / \code{\LinkA{load\_tolerance\_update()}{load.Rul.tolerance.Rul.update}}: Persist across sessions.
\item{} \code{\LinkA{reset\_tolerance\_update()}{reset.Rul.tolerance.Rul.update}}: Revert to built-in prior parameters.

\end{itemize}

\end{Section}
%
\begin{Section}{Risk and disturbance modules}

Optional scored overlays (all \code{fetch\_live = FALSE} by default):
\begin{itemize}

\item{} \code{\LinkA{score\_predation\_risk()}{score.Rul.predation.Rul.risk}}: Starfish, crab, snail predation pressure.
\item{} \code{\LinkA{score\_hab\_risk()}{score.Rul.hab.Rul.risk}}: Harmful algal bloom risk (PSP/ASP/DSP/AZP).
\item{} \code{\LinkA{score\_anthropogenic\_disturbance()}{score.Rul.anthropogenic.Rul.disturbance}}: Bottom trawling (ICES VMS SAR),
anchor damage, dredging/extraction.
\item{} \code{\LinkA{score\_wave\_exposure()}{score.Rul.wave.Rul.exposure}}: JONSWAP wave height from fetch and wind speed.
\item{} \code{\LinkA{score\_sediment\_stability()}{score.Rul.sediment.Rul.stability}}: Shields parameter mobility analysis.
\item{} \code{\LinkA{add\_shellfish\_classification()}{add.Rul.shellfish.Rul.classification}}: UK/EU harvesting area classification
(Class A/B/C/Prohibited; EC Regulation 854/2004).

\end{itemize}

\end{Section}
%
\begin{Section}{Live data integration}

Set credentials once; all live fetches are off by default:

\begin{alltt}oystermapR_live_config(cmems_user = "...", cmems_password = "...")
result_live <- fetch_live_environmental_data(survey)
\end{alltt}

\end{Section}
%
\begin{Section}{Multi-species comparison}

\begin{itemize}

\item{} \code{\LinkA{compare\_species()}{compare.Rul.species}}: Compare suitability across multiple species at the
same locations.
\item{} \code{\LinkA{compare\_surveys()}{compare.Rul.surveys}}: Compare two survey datasets for the same species.
\item{} \code{\LinkA{composite\_seasonal()}{composite.Rul.seasonal}}: Merge summer/winter surveys into a composite score.

\end{itemize}

\end{Section}
%
\begin{Author}
\strong{Maintainer}: T Tucker \email{tristantucker48@gmail.com}

\end{Author}
%
\begin{References}
Roberts D.R. et al. (2017) Cross-validation strategies for data with
temporal, spatial, hierarchical, or phylogenetic structure. Ecography
40:913-929. \Rhref{https://doi.org/10.1111/ecog.02881}{doi:10.1111\slash{}ecog.02881}

Breiman L. (2001) Random Forests. Machine Learning 45:5-32.

Gelman A. et al. (2013) Bayesian Data Analysis, 3rd ed. CRC Press.

Eigaard O.R. et al. (2017) The footprint of bottom trawling in European
waters. ICES J Mar Sci 74:1700-1710.

Hasselmann K. et al. (1973) Measurements of wind-wave growth and swell
decay during the Joint North Sea Wave Project (JONSWAP). Dtsch
Hydrogr Z Suppl A8.
\end{References}
\HeaderA{.apply\_bayesian\_update}{Apply any cached Bayesian updates to a tolerance list}{.apply.Rul.bayesian.Rul.update}
\keyword{internal}{.apply\_bayesian\_update}
%
\begin{Description}
Called internally by \code{\LinkA{get\_species\_tolerances()}{get.Rul.species.Rul.tolerances}} and \code{\LinkA{predict\_oyster()}{predict.Rul.oyster}} to
transparently substitute updated parameters when a Bayesian update exists
for a species in the session cache. End users rarely need to call this
directly.
\end{Description}
%
\begin{Usage}
\begin{verbatim}
.apply_bayesian_update(tols, species)
\end{verbatim}
\end{Usage}
%
\begin{Arguments}
\begin{ldescription}
\item[\code{tols}] Tolerance list from \code{\LinkA{get\_species\_tolerances()}{get.Rul.species.Rul.tolerances}}.

\item[\code{species}] Character. Species key.
\end{ldescription}
\end{Arguments}
%
\begin{Value}
Updated tolerance list (or unchanged if no cache entry found).
\end{Value}
\HeaderA{.ascii\_roc}{Print a simple ASCII ROC curve to the console}{.ascii.Rul.roc}
\keyword{internal}{.ascii\_roc}
%
\begin{Description}
Print a simple ASCII ROC curve to the console
\end{Description}
%
\begin{Usage}
\begin{verbatim}
.ascii_roc(fpr, tpr, auc, width = 50, height = 20)
\end{verbatim}
\end{Usage}
\HeaderA{.auto\_map\_columns}{Automatic column name matching for common sensor exports}{.auto.Rul.map.Rul.columns}
\keyword{internal}{.auto\_map\_columns}
%
\begin{Description}
Automatic column name matching for common sensor exports
\end{Description}
%
\begin{Usage}
\begin{verbatim}
.auto_map_columns(df, verbose = TRUE)
\end{verbatim}
\end{Usage}
\HeaderA{.detect\_sidelobe\_bins}{Internal: detect sidelobe-contaminated bins}{.detect.Rul.sidelobe.Rul.bins}
\keyword{internal}{.detect\_sidelobe\_bins}
%
\begin{Description}
Internal: detect sidelobe-contaminated bins
\end{Description}
%
\begin{Usage}
\begin{verbatim}
.detect_sidelobe_bins(speed_matrix, max_plausible_speed, sidelobe_threshold)
\end{verbatim}
\end{Usage}
%
\begin{Arguments}
\begin{ldescription}
\item[\code{speed\_matrix}] Numeric matrix, rows = ensembles, cols = bins (shallow to deep).

\item[\code{max\_plausible\_speed}] Numeric. Speed threshold in m/s.

\item[\code{sidelobe\_threshold}] Numeric. Contamination fraction threshold \LinkA{0,1}{0,1}.
\end{ldescription}
\end{Arguments}
%
\begin{Value}
Logical vector, length = n\_bins. TRUE = contaminated.
\end{Value}
\HeaderA{.distance\_to\_nearest\_m}{Internal: distance (metres) from each raster cell to the nearest survey point}{.distance.Rul.to.Rul.nearest.Rul.m}
\keyword{internal}{.distance\_to\_nearest\_m}
%
\begin{Description}
For each cell centre, finds the minimum Euclidean distance to any survey
point and converts to approximate metres. Used as the uncertainty band.
Chunked to stay within memory budget.
\end{Description}
%
\begin{Usage}
\begin{verbatim}
.distance_to_nearest_m(pts_lon, pts_lat, rast)
\end{verbatim}
\end{Usage}
\HeaderA{.find\_col\_any}{Find first matching column (case-insensitive)}{.find.Rul.col.Rul.any}
\keyword{internal}{.find\_col\_any}
%
\begin{Description}
Find first matching column (case-insensitive)
\end{Description}
%
\begin{Usage}
\begin{verbatim}
.find_col_any(df, candidates)
\end{verbatim}
\end{Usage}
\HeaderA{.get\_cached\_tolerances}{Retrieve cached tolerance list (or NULL if not updated)}{.get.Rul.cached.Rul.tolerances}
\keyword{internal}{.get\_cached\_tolerances}
%
\begin{Description}
Retrieve cached tolerance list (or NULL if not updated)
\end{Description}
%
\begin{Usage}
\begin{verbatim}
.get_cached_tolerances(species)
\end{verbatim}
\end{Usage}
\HeaderA{.haversine\_km}{Haversine distance in km between two WGS84 points}{.haversine.Rul.km}
\keyword{internal}{.haversine\_km}
%
\begin{Description}
Haversine distance in km between two WGS84 points
\end{Description}
%
\begin{Usage}
\begin{verbatim}
.haversine_km(lat1, lon1, lat2, lon2)
\end{verbatim}
\end{Usage}
\HeaderA{.idw\_interpolate}{Internal IDW interpolation (no external dependencies)}{.idw.Rul.interpolate}
\keyword{internal}{.idw\_interpolate}
%
\begin{Description}
Vectorised Inverse Distance Weighting using base-R matrix operations.
Computes a weighted mean of point values for every cell centre in a raster
template, using only points within \code{max\_dist} degrees of each cell.
\end{Description}
%
\begin{Usage}
\begin{verbatim}
.idw_interpolate(pts_lon, pts_lat, values, rast, power = 2, max_dist = 0.02)
\end{verbatim}
\end{Usage}
%
\begin{Arguments}
\begin{ldescription}
\item[\code{pts\_lon}, \code{pts\_lat}] Numeric vectors of point coordinates.

\item[\code{values}] Numeric vector of values to interpolate (same length as pts\_*).

\item[\code{rast}] A \code{SpatRaster} template defining the output grid.

\item[\code{power}] Numeric. Distance decay exponent.

\item[\code{max\_dist}] Numeric. Search radius in same units as coordinates.
\end{ldescription}
\end{Arguments}
%
\begin{Value}
Numeric vector of interpolated values, one per raster cell.
\end{Value}
\HeaderA{.mode}{Mode helper (most common value in a vector)}{.mode}
\keyword{internal}{.mode}
%
\begin{Description}
Mode helper (most common value in a vector)
\end{Description}
%
\begin{Usage}
\begin{verbatim}
.mode(x)
\end{verbatim}
\end{Usage}
\HeaderA{.nodal\_factors}{Compute Schureman nodal amplitude (f) and phase (u) corrections}{.nodal.Rul.factors}
\keyword{internal}{.nodal\_factors}
%
\begin{Description}
Computes the 18.6-year lunar nodal modulation factors f (amplitude scale)
and u (phase correction in degrees) for M2, S2, N2, K1, O1 at a given epoch.

Formulae from Schureman (1940) Tables 2 \& 3; same as those used by
UKHO, NOAA, and most modern tidal prediction software.
\end{Description}
%
\begin{Usage}
\begin{verbatim}
.nodal_factors(t_mean)
\end{verbatim}
\end{Usage}
%
\begin{Arguments}
\begin{ldescription}
\item[\code{t\_mean}] POSIXct. Representative timestamp for the survey (typically
the midpoint of the survey period). Nodal factors vary over months, so
a single midpoint evaluation is accurate to well within the constituent
uncertainty budget.
\end{ldescription}
\end{Arguments}
%
\begin{Value}
Named list with elements \code{f} and \code{u}, each a named numeric vector
indexed by constituent name (M2, S2, N2, K1, O1).
\end{Value}
\HeaderA{.numerical\_hessian}{Compute numerical Hessian by finite differences (central differences)}{.numerical.Rul.hessian}
\keyword{internal}{.numerical\_hessian}
%
\begin{Description}
Compute numerical Hessian by finite differences (central differences)
\end{Description}
%
\begin{Usage}
\begin{verbatim}
.numerical_hessian(f, theta, h = 1e-04)
\end{verbatim}
\end{Usage}
\HeaderA{.predict\_harmonic\_heights}{Predict tidal heights using embedded harmonic constituents (with nodal corrections)}{.predict.Rul.harmonic.Rul.heights}
\keyword{internal}{.predict\_harmonic\_heights}
%
\begin{Description}
Predict tidal heights using embedded harmonic constituents (with nodal corrections)
\end{Description}
%
\begin{Usage}
\begin{verbatim}
.predict_harmonic_heights(t_posix, port)
\end{verbatim}
\end{Usage}
%
\begin{Arguments}
\begin{ldescription}
\item[\code{t\_posix}] POSIXct vector of survey timestamps (UTC).

\item[\code{port}] Named list from .harmonic\_ports.
\end{ldescription}
\end{Arguments}
%
\begin{Value}
Numeric vector of predicted heights above Chart Datum (metres).
\end{Value}
\HeaderA{.rwmh}{Random-Walk Metropolis-Hastings MCMC sampler}{.rwmh}
\keyword{internal}{.rwmh}
%
\begin{Description}
Random-Walk Metropolis-Hastings MCMC sampler
\end{Description}
%
\begin{Usage}
\begin{verbatim}
.rwmh(
  log_post,
  theta_init,
  n_iter,
  n_warmup,
  proposal_sd,
  lower,
  upper,
  verbose = TRUE
)
\end{verbatim}
\end{Usage}
\HeaderA{.score\_categorical}{Score a categorical variable}{.score.Rul.categorical}
\keyword{internal}{.score\_categorical}
%
\begin{Description}
Score a categorical variable
\end{Description}
%
\begin{Usage}
\begin{verbatim}
.score_categorical(x, params)
\end{verbatim}
\end{Usage}
\HeaderA{.score\_numeric}{Score a single numeric variable against an optimal-range definition}{.score.Rul.numeric}
\keyword{internal}{.score\_numeric}
%
\begin{Description}
Score a single numeric variable against an optimal-range definition
\end{Description}
%
\begin{Usage}
\begin{verbatim}
.score_numeric(x, params)
\end{verbatim}
\end{Usage}
%
\begin{Arguments}
\begin{ldescription}
\item[\code{x}] Numeric value to score.

\item[\code{params}] Named list of scoring parameters.
\end{ldescription}
\end{Arguments}
%
\begin{Value}
Numeric in \LinkA{0, 1}{0, 1}. Returns \code{NA} if \code{x} is \code{NA}.
\end{Value}
\HeaderA{.score\_row}{Score all weighted factors for a single row (season-aware)}{.score.Rul.row}
\keyword{internal}{.score\_row}
%
\begin{Description}
Computes a weighted mean suitability score. Variables of type \code{"seasonal"}
have their parameters swapped from \code{tolerances\$seasonal\_overrides} when a
\code{season} value is present in the row. Missing columns are skipped and their
weight is excluded from the denominator.
\end{Description}
%
\begin{Usage}
\begin{verbatim}
.score_row(row, tolerances)
\end{verbatim}
\end{Usage}
%
\begin{Arguments}
\begin{ldescription}
\item[\code{row}] Named list or single-row dataframe slice.

\item[\code{tolerances}] Species tolerance list from \code{\LinkA{get\_species\_tolerances()}{get.Rul.species.Rul.tolerances}}.
\end{ldescription}
\end{Arguments}
%
\begin{Value}
List: \code{score}, \code{variable\_scores}, \code{variable\_weights}, \code{limiting\_factors}.
\end{Value}
\HeaderA{.species\_tolerances}{Species Tolerance Data for oystermapR}{.species.Rul.tolerances}
\keyword{internal}{.species\_tolerances}
%
\begin{Description}
Internal lookup tables defining exclusion thresholds, weighted scoring
parameters, and seasonal scoring overrides for each supported oyster species.

\strong{Structure per species entry:}
\begin{itemize}

\item{} \code{exclusions}         \bsl{}u2014 hard-stop limits; location is excluded if breached
\item{} \code{scored}             \bsl{}u2014 weighted variables contributing to the 0-1 score
\item{} \code{seasonal\_overrides} \bsl{}u2014 season-specific parameter swaps for scored variables

\end{itemize}


\strong{Adding a new species:} append a named list entry following the same
structure. Use \code{"seasonal\_overrides" = list()} if no seasonal data exists.
\end{Description}
%
\begin{Usage}
\begin{verbatim}
.species_tolerances
\end{verbatim}
\end{Usage}
%
\begin{Format}
A named list, one entry per species keyed by lowercase Latin name
with underscores (e.g. \code{"ostrea\_edulis"}).
\end{Format}
\HeaderA{.standardise\_coords}{Standardise coordinate column names to lat/lon}{.standardise.Rul.coords}
\keyword{internal}{.standardise\_coords}
%
\begin{Description}
Standardise coordinate column names to lat/lon
\end{Description}
%
\begin{Usage}
\begin{verbatim}
.standardise_coords(df)
\end{verbatim}
\end{Usage}
\HeaderA{.top\_introduction\_sites}{Find top N spatially distinct introduction sites with patch radius}{.top.Rul.introduction.Rul.sites}
\keyword{internal}{.top\_introduction\_sites}
%
\begin{Description}
Identifies the highest-scoring, spatially spread-out locations for oyster
introduction. Uses a greedy selection that prevents clustering \bsl{}u2014 once a site
is chosen, no other site within \code{min\_spacing\_deg} is considered, so the
result always represents distinct areas of the survey.

Patch radius is the equivalent circular radius of all suitable cells
(suitability >= \code{min\_suitability}) within the search neighbourhood. This
gives a practical sense of how large the suitable area is around each site.
\end{Description}
%
\begin{Usage}
\begin{verbatim}
.top_introduction_sites(
  df,
  n = 5L,
  min_suitability = 0.4,
  min_spacing_deg = 0.002,
  patch_search_deg = 0.005,
  spatial_res = 4L
)
\end{verbatim}
\end{Usage}
%
\begin{Arguments}
\begin{ldescription}
\item[\code{df}] Scored dataframe from predict\_oyster().

\item[\code{n}] Integer. Number of sites to return (default 5).

\item[\code{min\_suitability}] Numeric. Minimum score to consider suitable \LinkA{0,1}{0,1}.

\item[\code{min\_spacing\_deg}] Numeric. Minimum separation between selected sites
in decimal degrees (default 0.002 \bsl{}u2248 220 m). Prevents adjacent cells all
being reported as separate "sites".

\item[\code{patch\_search\_deg}] Numeric. Radius in degrees to count nearby suitable
cells when computing patch size (default 0.005 \bsl{}u2248 550 m).

\item[\code{spatial\_res}] Integer. Survey grid resolution in decimal places (default 4).
\end{ldescription}
\end{Arguments}
%
\begin{Value}
Dataframe of top sites, or NULL if none qualify.
\end{Value}
\HeaderA{add\_intertidal\_flag}{Add an intertidal zone flag to a survey dataframe}{add.Rul.intertidal.Rul.flag}
%
\begin{Description}
Classifies each survey location as intertidal, subtidal, or supratidal
based on its depth below Chart Datum (CD) and the tidal range at the
nearest harmonic reference port. Chart Datum is approximately Lowest
Astronomical Tide (LAT), so:
\begin{itemize}

\item{} \strong{Subtidal:}    \code{depth > 0} (always submerged)
\item{} \strong{Intertidal:}  \AsIs{\texttt{-MHWS\_above\_CD <= depth <= 0}}
(exposed at some states of the tide; depth is at or above CD)
\item{} \strong{Supratidal:}  \code{depth < -MHWS\_above\_CD}
(above the highest astronomical tide; excluded from scoring)

\end{itemize}


MHWS above CD is approximated from the harmonic constituents of the
nearest port as: \code{Z0 + 1.1 * (M2\_H + S2\_H)} \bsl{}u2014 the 10\% factor accounts
for minor constituents not included in the 5-constituent model.

This flag is used internally by \code{\LinkA{score\_locations()}{score.Rul.locations}} to give \emph{Magallana
gigas} (Pacific oyster) full depth scores for intertidal cells, reflecting
its strong intertidal ecology in NW European waters.
\end{Description}
%
\begin{Usage}
\begin{verbatim}
add_intertidal_flag(
  df,
  max_port_dist_km = 75,
  depth_col = "depth",
  verbose = TRUE
)
\end{verbatim}
\end{Usage}
%
\begin{Arguments}
\begin{ldescription}
\item[\code{df}] Dataframe with \code{lat}, \code{lon}, and \code{depth} (chart-datum corrected)
columns. Typically the output of \code{\LinkA{auto\_tidal\_correct()}{auto.Rul.tidal.Rul.correct}} or
\code{\LinkA{correct\_to\_chart\_datum()}{correct.Rul.to.Rul.chart.Rul.datum}}.

\item[\code{max\_port\_dist\_km}] Numeric. Maximum distance to nearest harmonic port
in km (default 75). If the survey centroid is further than this from all
ports, the function falls back to a conservative intertidal window of 6 m
above CD and issues a warning.

\item[\code{depth\_col}] Character. Name of the chart-datum depth column
(default \code{"depth"}).

\item[\code{verbose}] Logical. Print the nearest port name and derived MHWS
(default \code{TRUE}).
\end{ldescription}
\end{Arguments}
%
\begin{Value}
The input dataframe with an added \code{intertidal\_zone} column:
\begin{itemize}

\item{} \code{"subtidal"}: depth > 0 m CD
\item{} \code{"intertidal"}: 0 m CD >= depth >= -MHWS\_above\_CD
\item{} \code{"supratidal"}: depth < -MHWS\_above\_CD

\end{itemize}

\end{Value}
%
\begin{Examples}
\begin{ExampleCode}
## Not run: 
survey <- auto_tidal_correct(survey, datetime_col = "date")
survey <- add_intertidal_flag(survey)

# Check intertidal coverage
table(survey$intertidal_zone)

## End(Not run)
\end{ExampleCode}
\end{Examples}
\HeaderA{add\_season\_column}{Detect season for each row in a dataframe}{add.Rul.season.Rul.column}
%
\begin{Description}
Vectorised wrapper around \code{\LinkA{detect\_season()}{detect.Rul.season}} for use on dataframes.
\end{Description}
%
\begin{Usage}
\begin{verbatim}
add_season_column(df, date_col = NULL, lat_col = "lat")
\end{verbatim}
\end{Usage}
%
\begin{Arguments}
\begin{ldescription}
\item[\code{df}] A dataframe containing at minimum a date column and a latitude column.

\item[\code{date\_col}] Name of the date column (default \code{"date"}).

\item[\code{lat\_col}] Name of the latitude column (default \code{"lat"}).
\end{ldescription}
\end{Arguments}
%
\begin{Value}
The input dataframe with an additional column \code{season}.
\end{Value}
\HeaderA{add\_shellfish\_classification}{Add shellfish water quality classification to scored result}{add.Rul.shellfish.Rul.classification}
%
\begin{Description}
Attaches a regulatory shellfish harvesting classification (A/B/C/Prohibited)
to each row of a scored result dataframe. This classification affects the
economic viability score for aquaculture applications: Prohibited sites
receive viability = 0; Class C sites receive a 0.60 multiplier; Class B a
0.80 multiplier; Class A is unpenalised.

Classification can be supplied three ways (in priority order):
\begin{enumerate}

\item{} \strong{\code{class\_col}} \bsl{}u2014 name of an existing column in \code{result} already
containing classification values ("A", "B", "C", "Prohibited", or NA).
\item{} \strong{\code{classified\_areas}} \bsl{}u2014 a dataframe with \code{lat}, \code{lon}, and
\code{shellfish\_class} columns, spatially matched to result rows within
\code{match\_radius\_deg} degrees.
\item{} \strong{\code{fetch\_live = TRUE}} \bsl{}u2014 queries the FSA England/Wales open data API
(requires internet). Requires \code{httr} package. Works only for sites
within the FSA/DAERA coverage area (England, Wales, Northern Ireland).

\end{enumerate}


Unclassified or unmatched sites receive \code{shellfish\_class = "Unclassified"}
and a penalty of 0.70 (precautionary principle).
\end{Description}
%
\begin{Usage}
\begin{verbatim}
add_shellfish_classification(
  result,
  class_col = NULL,
  classified_areas = NULL,
  fetch_live = FALSE,
  match_radius_deg = 0.05,
  verbose = TRUE
)
\end{verbatim}
\end{Usage}
%
\begin{Arguments}
\begin{ldescription}
\item[\code{result}] Dataframe from \code{\LinkA{predict\_oyster()}{predict.Rul.oyster}} or \code{\LinkA{score\_locations()}{score.Rul.locations}}.
Must contain \code{lat} and \code{lon}.

\item[\code{class\_col}] Character. Name of an existing column in \code{result} containing
class values. If supplied, \code{classified\_areas} and \code{fetch\_live} are ignored.

\item[\code{classified\_areas}] Dataframe with \code{lat}, \code{lon}, \code{shellfish\_class}
columns. Used when \code{class\_col} is NULL and \code{fetch\_live = FALSE}.

\item[\code{fetch\_live}] Logical. Query live FSA/DAERA API (default FALSE). Requires
internet access and the \code{httr} package. If fetch fails, falls back to
"Unclassified".

\item[\code{match\_radius\_deg}] Numeric. Spatial matching tolerance in decimal
degrees when using \code{classified\_areas} (default 0.05 = \textasciitilde{}5 km).

\item[\code{verbose}] Logical. Print matching summary (default TRUE).
\end{ldescription}
\end{Arguments}
%
\begin{Value}
\code{result} with two additional columns:
\begin{itemize}

\item{} \code{shellfish\_class}: "A", "B", "C", "Prohibited", or "Unclassified"
\item{} \code{shellfish\_class\_penalty}: multiplier applied to economic viability
(1.0 for A, 0.80 for B, 0.60 for C, 0.0 for Prohibited, 0.70 for
Unclassified)

\end{itemize}

\end{Value}
%
\begin{Examples}
\begin{ExampleCode}
## Not run: 
# Option 1: manual column already in data
result <- add_shellfish_classification(result, class_col = "water_class")

# Option 2: separate classified areas dataframe
areas <- data.frame(lat = c(51.5), lon = c(-4.2), shellfish_class = c("B"))
result <- add_shellfish_classification(result, classified_areas = areas)

# Option 3: live fetch (internet required)
result <- add_shellfish_classification(result, fetch_live = TRUE)

## End(Not run)
\end{ExampleCode}
\end{Examples}
\HeaderA{add\_suitability\_ci}{Add bootstrap confidence intervals to suitability scores}{add.Rul.suitability.Rul.ci}
%
\begin{Description}
Re-scores each row of a survey dataframe \code{n\_boot} times, each time adding
Gaussian noise to every numeric variable proportional to its measurement
uncertainty, and returns the 5th and 95th percentile suitability scores as
\code{suit\_ci\_lower} and \code{suit\_ci\_upper} columns (90\% interval).

Measurement uncertainty is specified via the \code{uncertainty} argument: a named
list where each name matches a variable name in \code{tolerances\$scored\_factors}
and the value is the 1-sigma absolute error (same units as the variable).
Any variable not listed defaults to 2\% of its observed value.
\end{Description}
%
\begin{Usage}
\begin{verbatim}
add_suitability_ci(
  result,
  tolerances,
  n_boot = 200L,
  uncertainty = NULL,
  seed = 42L,
  verbose = FALSE
)
\end{verbatim}
\end{Usage}
%
\begin{Arguments}
\begin{ldescription}
\item[\code{result}] Dataframe returned by \code{\LinkA{predict\_oyster()}{predict.Rul.oyster}}.

\item[\code{tolerances}] Species tolerance list from \code{\LinkA{get\_species\_tolerances()}{get.Rul.species.Rul.tolerances}}.

\item[\code{n\_boot}] Integer. Number of bootstrap resamples (default 200).

\item[\code{uncertainty}] Named numeric vector. 1-sigma measurement errors by
variable name (e.g. \code{c(temperature = 0.2, salinity = 0.5)}).

\item[\code{seed}] Integer or NULL. Random seed for reproducibility (default 42).

\item[\code{verbose}] Logical. Print progress (default FALSE).
\end{ldescription}
\end{Arguments}
%
\begin{Value}
The input dataframe with additional columns:
\code{suit\_ci\_lower}, \code{suit\_ci\_upper}, \code{suit\_ci\_width}.
\end{Value}
%
\begin{Examples}
\begin{ExampleCode}
## Not run: 
result <- predict_oyster(survey, "ostrea_edulis")
tol    <- get_species_tolerances("ostrea_edulis")
result <- add_suitability_ci(result, tol,
                             uncertainty = c(temperature = 0.3,
                                             salinity    = 0.5))
# Columns suit_ci_lower, suit_ci_upper, suit_ci_width now present
generate_report(result, "report.html")   # map rings scale with CI width

## End(Not run)
\end{ExampleCode}
\end{Examples}
\HeaderA{analyse\_connectivity}{Analyse spatial connectivity of suitable habitat cells}{analyse.Rul.connectivity}
%
\begin{Description}
Builds a spatial graph of cells above a suitability threshold and identifies
connected components \bsl{}u2014 contiguous patches of suitable habitat. Isolated
cells or small patches are flagged because they are at higher risk of:
\begin{itemize}

\item{} \strong{Recruitment failure} \bsl{}u2014 larvae from isolated populations may be
flushed away before settling back in the patch.
\item{} \strong{Genetic bottlenecks} \bsl{}u2014 small isolated populations have lower allelic
diversity and reduced adaptive capacity.
\item{} \strong{Local extinction without recolonisation} \bsl{}u2014 if a patch is lost, there
are no connected source populations to recolonise.

\end{itemize}


Two cells are considered connected if they are within \code{gap\_m} metres of
each other (default 500 m \bsl{}u2014 approximately the scale of larval dispersal
during a tidal cycle at moderate current speeds). Connectivity is computed
using a fast union-find (disjoint set) algorithm \bsl{}u2014 no external graph
packages required.
\end{Description}
%
\begin{Usage}
\begin{verbatim}
analyse_connectivity(
  result,
  min_suitability = 0.5,
  gap_m = 500,
  min_patch_cells = 3L,
  verbose = TRUE
)
\end{verbatim}
\end{Usage}
%
\begin{Arguments}
\begin{ldescription}
\item[\code{result}] Dataframe from \code{\LinkA{predict\_oyster()}{predict.Rul.oyster}} with \code{lat}, \code{lon},
\code{suitability}, and \code{suitability\_class} columns.

\item[\code{min\_suitability}] Numeric. Minimum suitability score for a cell to be
included in the connectivity analysis (default 0.5, i.e. Moderate+).

\item[\code{gap\_m}] Numeric. Maximum distance in metres between two suitable cells
for them to be considered connected (default 500 m).

\item[\code{min\_patch\_cells}] Integer. Patches smaller than this are flagged as
isolated (default 3 cells).

\item[\code{verbose}] Logical. Print connectivity summary (default TRUE).
\end{ldescription}
\end{Arguments}
%
\begin{Value}
Input dataframe with additional columns:
\code{patch\_id} (integer \bsl{}u2014 unique patch identifier; NA for non-suitable cells),
\code{patch\_size} (number of cells in the patch),
\code{patch\_area\_km2} (approximate area based on cell density),
\code{connectivity\_class} (\code{"isolated"}, \code{"small"}, \code{"moderate"}, \code{"large"}),
\code{is\_hub} (logical \bsl{}u2014 TRUE for the highest-scoring cell in each patch,
useful as candidate introduction points).
\end{Value}
%
\begin{Examples}
\begin{ExampleCode}
## Not run: 
result <- predict_oyster(survey, "ostrea_edulis")
result <- analyse_connectivity(result, gap_m = 500)

# Large well-connected patches are the best restoration targets
good_patches <- subset(result,
  connectivity_class == "large" & suitability_class == "High")

# Count patches
table(result$connectivity_class)

# Visualise \u2014 patch_id maps to distinct colours in QGIS
export_geotiff(result, "connectivity.tif")

## End(Not run)
\end{ExampleCode}
\end{Examples}
\HeaderA{assess\_gear\_feasibility}{Assess gear deployment feasibility at survey locations}{assess.Rul.gear.Rul.feasibility}
%
\begin{Description}
Evaluates which aquaculture and restoration gear types are physically viable
at each survey location based on depth, current velocity, substrate, and
(optionally) wave height. Returns a feasibility score and binary flag for
each gear type.

This function is part of the \strong{planning module} aimed at commercial
operators. For pure restoration work, the \code{restoration\_reef} gear type is
always included and is the most relevant output.
\end{Description}
%
\begin{Usage}
\begin{verbatim}
assess_gear_feasibility(
  result,
  gear_types = names(.gear_specs),
  depth_col = "depth",
  current_col = "current_velocity",
  substrate_col = "substrate_hardness_class",
  wave_col = NULL,
  verbose = TRUE
)
\end{verbatim}
\end{Usage}
%
\begin{Arguments}
\begin{ldescription}
\item[\code{result}] Dataframe from \code{\LinkA{predict\_oyster()}{predict.Rul.oyster}} with environmental columns.

\item[\code{gear\_types}] Character vector of gear types to assess. Options:
\code{"longline"}, \code{"bottom\_culture"}, \code{"intertidal\_rack"}, \code{"raft"},
\code{"restoration\_reef"}. Default: all five.

\item[\code{depth\_col}] Character. Column name for depth below CD in metres
(default \code{"depth"}; positive = subtidal).

\item[\code{current\_col}] Character. Column name for current velocity in m/s
(default \code{"current\_velocity"}).

\item[\code{substrate\_col}] Character. Column for substrate class
(default \code{"substrate\_hardness\_class"}). Can be the output of
\code{\LinkA{classify\_substrate\_from\_backscatter()}{classify.Rul.substrate.Rul.from.Rul.backscatter}}.

\item[\code{wave\_col}] Character or NULL. Column for significant wave height Hs
in metres (default NULL \bsl{}u2014 wave check skipped if absent).

\item[\code{verbose}] Logical. Print feasibility summary (default TRUE).
\end{ldescription}
\end{Arguments}
%
\begin{Value}
Input dataframe with added columns per gear type:
\AsIs{\texttt{gear\_<type>\_feasible}} (logical),
\AsIs{\texttt{gear\_<type>\_score}} \LinkA{0-1}{0.Rdash.1},
plus \code{best\_gear} (most feasible option) and \code{n\_gear\_options} (count).
\end{Value}
%
\begin{Examples}
\begin{ExampleCode}
## Not run: 
result <- predict_oyster(survey, "ostrea_edulis")
result <- assess_gear_feasibility(result)

# Sites where longline is feasible AND suitability is High
longline_sites <- subset(result,
  gear_longline_feasible & suitability_class == "High")

# Restoration reef sites (no farming licence needed)
reef_sites <- subset(result, gear_restoration_reef_feasible)

## End(Not run)
\end{ExampleCode}
\end{Examples}
\HeaderA{auto\_tidal\_correct}{Automatically predict and apply tidal correction from harmonic constituents}{auto.Rul.tidal.Rul.correct}
%
\begin{Description}
Finds the nearest standard port to the survey area, checks whether it is
within a configurable distance threshold, predicts tidal heights for every
survey timestamp using embedded 5-constituent harmonic models
(M2, S2, N2, K1, O1), and applies \code{\LinkA{correct\_to\_chart\_datum()}{correct.Rul.to.Rul.chart.Rul.datum}} automatically.

No external data or internet connection is required. Harmonic constituent
data (amplitude and Greenwich phase lag) for 31 UK and European ports are
embedded in the package.

\strong{Accuracy:} 5-constituent predictions are typically accurate to \bsl{}u00b10.15\bsl{}u20130.35 m
at the standard port. Accuracy degrades with distance from the port and in
areas with complex tidal dynamics (e.g. the Bristol Channel, inner fjords).
For safety-critical work, use official UKHO EasyTide predictions and supply
them manually via \code{\LinkA{correct\_to\_chart\_datum()}{correct.Rul.to.Rul.chart.Rul.datum}}.

\strong{Requirements:} The dataframe must contain a datetime column. Values must
be parseable by \code{\LinkA{as.POSIXct()}{as.POSIXct}} (e.g. \code{"2024-06-15 10:30:00"} or
\code{"2024-06-15"}). Date-only strings default to noon UTC.
\end{Description}
%
\begin{Usage}
\begin{verbatim}
auto_tidal_correct(
  df,
  datetime_col = "date",
  depth_col = "depth",
  max_port_dist_km = 75,
  keep_raw = TRUE,
  verbose = TRUE
)
\end{verbatim}
\end{Usage}
%
\begin{Arguments}
\begin{ldescription}
\item[\code{df}] A dataframe with at minimum a depth and a datetime column.
Typically the output of \code{\LinkA{merge\_sensor\_data()}{merge.Rul.sensor.Rul.data}}.

\item[\code{datetime\_col}] Character. Name of the datetime column (default
\code{"date"}). Values should be UTC or treated as UTC.

\item[\code{depth\_col}] Character. Name of the depth column to correct
(default \code{"depth"}).

\item[\code{max\_port\_dist\_km}] Numeric. Maximum distance in km from the survey
centroid to the nearest standard port. If the nearest port is further
than this, no correction is applied and a warning is issued
(default \code{75} km).

\item[\code{keep\_raw}] Logical. Retain original depths in a \code{depth\_raw\_m} column
(default \code{TRUE}).

\item[\code{verbose}] Logical. Print port selection, distance, and prediction
summary (default \code{TRUE}).
\end{ldescription}
\end{Arguments}
%
\begin{Value}
The input dataframe with depths corrected to Chart Datum, plus
a \code{tidal\_height\_pred\_m} column containing the predicted tidal height
used for each row.
\end{Value}
%
\begin{SeeAlso}
\code{\LinkA{correct\_to\_chart\_datum()}{correct.Rul.to.Rul.chart.Rul.datum}} for manual correction,
\code{\LinkA{tidal\_port\_info()}{tidal.Rul.port.Rul.info}} for port datum reference data.
\end{SeeAlso}
%
\begin{Examples}
\begin{ExampleCode}
## Not run: 
# Automatic correction \u2014 finds nearest port, predicts heights, corrects depths
survey_corrected <- auto_tidal_correct(survey, datetime_col = "date")

# Tighten the distance threshold (only trust ports within 40 km)
survey_corrected <- auto_tidal_correct(survey, max_port_dist_km = 40)

# See which port was selected and inspect predicted heights
survey_corrected$tidal_height_pred_m

## End(Not run)
\end{ExampleCode}
\end{Examples}
\HeaderA{check\_exclusions}{Check exclusion criteria for all rows in a dataset}{check.Rul.exclusions}
%
\begin{Description}
Applies hard-stop exclusion criteria for a given species tolerance profile.
Any location that violates one or more exclusion rules receives a flag and
is removed from the suitability scoring pipeline.

Exclusion factors checked (where data columns are present):
\begin{itemize}

\item{} General temperature range
\item{} Season-specific temperature floors/ceilings (requires \code{season} column)
\item{} Salinity (with temperature-dependent minimum)
\item{} Dissolved oxygen (hypoxia threshold)

\end{itemize}

\end{Description}
%
\begin{Usage}
\begin{verbatim}
check_exclusions(df, tolerances)
\end{verbatim}
\end{Usage}
%
\begin{Arguments}
\begin{ldescription}
\item[\code{df}] A dataframe. Must contain at minimum a subset of the standard
column names (see \strong{Column Naming} below). Any exclusion factor for
which no column is present is silently skipped with a warning.

\item[\code{tolerances}] A species tolerance list as returned by
\code{\LinkA{get\_species\_tolerances()}{get.Rul.species.Rul.tolerances}}.
\end{ldescription}
\end{Arguments}
%
\begin{Value}
The input dataframe with two additional columns:
\begin{itemize}

\item{} \code{excluded}: logical \bsl{}u2014 \code{TRUE} if the location fails any exclusion criterion.
\item{} \code{exclusion\_reason}: character \bsl{}u2014 semicolon-separated list of failed criteria,
or \code{NA} if not excluded.

\end{itemize}

\end{Value}
%
\begin{Section}{Column Naming}

The function recognises these column names (case-insensitive):
\Tabular{ll}{
Variable & Expected column name(s) \\{}
Temperature & \code{temperature}, \code{temp} \\{}
Salinity & \code{salinity}, \code{sal} \\{}
Dissolved oxygen & \code{dissolved\_oxygen}, \code{do}, \code{do\_mgl} \\{}
Season & \code{season} (auto-added by \code{\LinkA{add\_season\_column()}{add.Rul.season.Rul.column}}) \\{}
}
\end{Section}
%
\begin{Examples}
\begin{ExampleCode}
tol <- get_species_tolerances("ostrea_edulis")
df <- data.frame(
  lat = 51.5, lon = -2.5,
  temperature = 8, salinity = 28, dissolved_oxygen = 7
)
check_exclusions(df, tol)
\end{ExampleCode}
\end{Examples}
\HeaderA{classify\_substrate\_from\_backscatter}{Classify seabed substrate hardness from near-seabed ADCP backscatter}{classify.Rul.substrate.Rul.from.Rul.backscatter}
%
\begin{Description}
Uses the intensity of acoustic backscatter from the near-seabed bin(s) to
classify substrate type. Hard substrates (gravel, rock, shell hash) scatter
more sound energy than soft substrates (mud, fine sand). The function
converts raw backscatter (dB re 1 m\textasciicircum{}-1 or instrument counts) to a
categorical substrate class and a continuous hardness index \LinkA{0, 1}{0, 1}.

\strong{Classification thresholds (default, adjustable):}
\Tabular{lll}{
Class & Hardness index & Typical substrate \\{}
Hard & 0.8 \bsl{}u2013 1.0 & Rock, cobble, shell hash \\{}
Mixed & 0.5 \bsl{}u2013 0.8 & Gravel/sand mix, maerl beds \\{}
Soft & 0.2 \bsl{}u2013 0.5 & Sand, sandy mud \\{}
Very Soft & 0.0 \bsl{}u2013 0.2 & Mud, fine silt \\{}
}


The hardness index is a min-max normalisation of the mean near-seabed
backscatter across the bottom \code{n\_bottom\_bins} cells. If the column
contains calibrated volume backscatter (Sv in dB), set \code{is\_sv = TRUE}
and a sign-flip is applied (less negative Sv = stronger return = harder).
\end{Description}
%
\begin{Usage}
\begin{verbatim}
classify_substrate_from_backscatter(
  df,
  backscatter_col,
  is_sv = FALSE,
  bs_min = NULL,
  bs_max = NULL,
  hard_thresh = 0.8,
  mixed_thresh = 0.5,
  soft_thresh = 0.2,
  output_col = "substrate_hardness",
  verbose = TRUE
)
\end{verbatim}
\end{Usage}
%
\begin{Arguments}
\begin{ldescription}
\item[\code{df}] Dataframe with at least one column of seabed backscatter values.

\item[\code{backscatter\_col}] Character. Name of the backscatter column.

\item[\code{is\_sv}] Logical. If TRUE, treats values as Sv (dB, typically negative);
hardness is derived from abs(Sv) after sign convention correction.
Default FALSE (raw instrument counts, higher = harder).

\item[\code{bs\_min}] Numeric. Backscatter value mapped to hardness = 0 (softest).
Defaults to 5th percentile of observed data.

\item[\code{bs\_max}] Numeric. Backscatter value mapped to hardness = 1 (hardest).
Defaults to 95th percentile of observed data.

\item[\code{hard\_thresh}] Numeric. Hardness index above which substrate is "Hard"
(default 0.8).

\item[\code{mixed\_thresh}] Numeric. Hardness index above which substrate is "Mixed"
(default 0.5).

\item[\code{soft\_thresh}] Numeric. Hardness index above which substrate is "Soft"
(default 0.2); below is "Very Soft".

\item[\code{output\_col}] Character. Base name for output columns. Two columns are
added: \AsIs{\texttt{<output\_col>\_index}} (numeric, 0-1) and \AsIs{\texttt{<output\_col>\_class}}
(character). Default \code{"substrate\_hardness"}.

\item[\code{verbose}] Logical. Print classification summary (default TRUE).
\end{ldescription}
\end{Arguments}
%
\begin{Value}
Input dataframe with \code{substrate\_hardness\_index} and
\code{substrate\_hardness\_class} columns appended. These map directly to the
\code{substrate\_hardness} scored variable in the tolerance spec.
\end{Value}
%
\begin{Examples}
\begin{ExampleCode}
## Not run: 
# ADCP data with near-seabed backscatter column
survey <- classify_substrate_from_backscatter(survey,
            backscatter_col = "bottom_backscatter_db",
            is_sv = TRUE)

# The substrate_hardness_index column feeds directly into predict_oyster()
# as the 'substrate_hardness' variable
names(survey)[names(survey) == "substrate_hardness_index"] <- "substrate_hardness"
result <- predict_oyster(survey, "ostrea_edulis")

## End(Not run)
\end{ExampleCode}
\end{Examples}
\HeaderA{compare\_species}{Compare suitability across multiple oyster species}{compare.Rul.species}
%
\begin{Description}
Runs \code{\LinkA{predict\_oyster()}{predict.Rul.oyster}} for every specified species against the same survey
dataset and returns a combined comparison table \bsl{}u2014 one row per survey
location, one suitability column per species. An optional summary prints
which species is best-suited to each location.

This is the tool to use when you haven't yet decided which species to stock,
or when you want to identify areas where multiple species overlap in
suitability (robust site selection).
\end{Description}
%
\begin{Usage}
\begin{verbatim}
compare_species(
  data,
  species = NULL,
  output_dir = NULL,
  min_data_quality = "low",
  verbose = TRUE
)
\end{verbatim}
\end{Usage}
%
\begin{Arguments}
\begin{ldescription}
\item[\code{data}] A dataframe or CSV path accepted by \code{\LinkA{predict\_oyster()}{predict.Rul.oyster}}.

\item[\code{species}] Character vector of species keys to compare. Defaults to all
species currently in the oystermapR database (see \code{\LinkA{list\_species()}{list.Rul.species}}).
Passing \code{NULL} uses all available species.

\item[\code{output\_dir}] Character or \code{NULL}. If a directory path is supplied,
individual GeoTIFF heatmaps are exported for each species into that folder.
Default \code{NULL} (no rasters written).

\item[\code{min\_data\_quality}] Character. Minimum data quality tier to include:
\code{"low"}, \code{"medium"}, or \code{"high"} (default \code{"low"} \bsl{}u2014 include all species).
Set to \code{"high"} to restrict to well-characterised species only.

\item[\code{verbose}] Logical. Print a comparison summary per species (default \code{TRUE}).
\end{ldescription}
\end{Arguments}
%
\begin{Value}
A dataframe with all original survey columns plus:
\begin{itemize}

\item{} \AsIs{\texttt{suit\_<species\_key>}}: suitability score \LinkA{0, 1}{0, 1} per species
\item{} \AsIs{\texttt{class\_<species\_key>}}: suitability class per species
\item{} \code{best\_species}: key of the species with the highest suitability at each
location (NA if all excluded)
\item{} \code{best\_suitability}: the highest suitability score across all species
\item{} \code{n\_species\_high}: number of species scoring "High" at each location
(useful for identifying robust multi-species sites)

\end{itemize}

\end{Value}
%
\begin{Examples}
\begin{ExampleCode}
## Not run: 
# Compare all supported species
comparison <- compare_species(survey)

# Compare only well-characterised species
comparison <- compare_species(survey, min_data_quality = "high")

# Compare specific species and export per-species heatmaps
comparison <- compare_species(
  survey,
  species    = c("ostrea_edulis", "magallana_gigas"),
  output_dir = "comparison_maps/"
)

# Find sites suitable for both O. edulis AND M. gigas
both_high <- subset(comparison, n_species_high >= 2)

## End(Not run)
\end{ExampleCode}
\end{Examples}
\HeaderA{compare\_surveys}{Compare suitability scores across multiple surveys or monitoring years}{compare.Rul.surveys}
%
\begin{Description}
Takes a named list of \code{predict\_oyster()} result dataframes and produces a
single comparative summary suitable for trend monitoring, multi-site
selection, or regulatory reporting. Common spatial cells are matched across
surveys and a change analysis is run between consecutive surveys if ordered
chronologically.

The function returns a list with three elements:
\begin{itemize}

\item{} \code{summary} \bsl{}u2014 one row per survey with mean suitability, class breakdown,
and data completeness stats.
\item{} \code{spatial} \bsl{}u2014 all surveys joined on shared lat/lon cells, with suitability
columns per survey and \code{trend\_slope} (linear regression of suitability
over survey index for each cell).
\item{} \code{change} \bsl{}u2014 pairwise change table (only if surveys are in order; each
consecutive pair shows mean score difference and fraction of cells
that improved, degraded, or were stable).

\end{itemize}

\end{Description}
%
\begin{Usage}
\begin{verbatim}
compare_surveys(
  surveys,
  cell_deg = 0.001,
  stable_threshold = 0.05,
  verbose = TRUE
)
\end{verbatim}
\end{Usage}
%
\begin{Arguments}
\begin{ldescription}
\item[\code{surveys}] Named list of dataframes from \code{\LinkA{predict\_oyster()}{predict.Rul.oyster}}. Names
should be meaningful labels (e.g. survey year or site name).

\item[\code{cell\_deg}] Numeric. Spatial rounding tolerance in decimal degrees for
matching cells across surveys (default 0.001 \bsl{}u2248 111 m).

\item[\code{stable\_threshold}] Numeric. Absolute suitability change below which a
cell is considered "stable" rather than improved/degraded (default 0.05).

\item[\code{verbose}] Logical. Print comparison summary (default TRUE).
\end{ldescription}
\end{Arguments}
%
\begin{Value}
Named list: \code{summary} (dataframe), \code{spatial} (dataframe),
\code{change} (dataframe or NULL if only one survey).
\end{Value}
%
\begin{Examples}
\begin{ExampleCode}
## Not run: 
r2022 <- predict_oyster(survey_2022, "ostrea_edulis")
r2023 <- predict_oyster(survey_2023, "ostrea_edulis")
r2024 <- predict_oyster(survey_2024, "ostrea_edulis")

comp <- compare_surveys(
  surveys = list("2022" = r2022, "2023" = r2023, "2024" = r2024)
)

comp$summary        # mean scores by year
comp$change         # year-on-year change

# Pass to generate_report() for a full comparative HTML report
generate_report(r2024, "monitoring_report.html",
                title = "Kames Bay 3-Year Monitoring")

## End(Not run)
\end{ExampleCode}
\end{Examples}
\HeaderA{composite\_seasonal}{Combine multi-season survey results into a composite suitability score}{composite.Rul.seasonal}
%
\begin{Description}
Oysters are permanent residents; a location is only truly suitable if
conditions are adequate year-round. \code{composite\_seasonal()} takes a named
list of \code{predict\_oyster()} result dataframes (one per season or survey
visit) and produces a single composite dataframe with one row per spatial
cell, summarising suitability across all input surveys.

\strong{Composite methods:}
\begin{itemize}

\item{} \code{"min"} (default) \bsl{}u2014 most conservative; a location must be suitable in
\emph{every} season to score well. Best for regulatory site selection.
\item{} \code{"mean"} \bsl{}u2014 average across seasons. Good for monitoring trend summaries.
\item{} \code{"prob"} \bsl{}u2014 fraction of seasons with suitability >= \code{prob\_threshold}.
Useful when surveys are irregularly spaced or some seasons are missing.

\end{itemize}


Spatial matching uses a grid cell tolerance of \code{cell\_deg} decimal degrees
(default 0.001 \bsl{}u2248 111 m). Cells not present in all seasons are flagged in
\code{n\_seasons\_present}.
\end{Description}
%
\begin{Usage}
\begin{verbatim}
composite_seasonal(
  surveys,
  method = c("min", "mean", "prob"),
  prob_threshold = 0.5,
  cell_deg = 0.001,
  verbose = TRUE
)
\end{verbatim}
\end{Usage}
%
\begin{Arguments}
\begin{ldescription}
\item[\code{surveys}] Named list of dataframes from \code{\LinkA{predict\_oyster()}{predict.Rul.oyster}}. Each must
have \code{lat}, \code{lon}, and \code{suitability} columns. Names are used as season
labels (e.g. \code{list(spring = r1, summer = r2, autumn = r3)}).

\item[\code{method}] Character. One of \code{"min"}, \code{"mean"}, \code{"prob"} (default \code{"min"}).

\item[\code{prob\_threshold}] Numeric \LinkA{0,1}{0,1}. Suitability threshold used when
\code{method = "prob"} (default 0.5).

\item[\code{cell\_deg}] Numeric. Spatial tolerance in decimal degrees for matching
cells across surveys (default 0.001).

\item[\code{verbose}] Logical. Print summary (default TRUE).
\end{ldescription}
\end{Arguments}
%
\begin{Value}
A dataframe with columns: \code{lat}, \code{lon}, \code{suitability\_composite},
\code{suitability\_class}, \code{n\_seasons\_present}, one \AsIs{\texttt{suit\_<season>}} column per
input survey, and \code{suit\_range} (max \bsl{}u2212 min across seasons).
\end{Value}
%
\begin{Examples}
\begin{ExampleCode}
## Not run: 
spring <- predict_oyster(survey_spring, "ostrea_edulis")
summer <- predict_oyster(survey_summer, "ostrea_edulis")
autumn <- predict_oyster(survey_autumn, "ostrea_edulis")

composite <- composite_seasonal(
  surveys = list(spring = spring, summer = summer, autumn = autumn),
  method  = "min"
)

generate_report(composite, "composite_report.html",
                title = "Year-round Suitability Assessment")

## End(Not run)
\end{ExampleCode}
\end{Examples}
\HeaderA{correct\_to\_chart\_datum}{Correct survey depths to Chart Datum (LAT)}{correct.Rul.to.Rul.chart.Rul.datum}
%
\begin{Description}
Converts water depths recorded during a survey to depths below Chart Datum
(approximately Lowest Astronomical Tide, LAT) by adding the tidal height
above Chart Datum at the time of survey. This removes bias introduced by
surveying at different states of the tide.

\strong{Formula:} \code{depth\_CD = depth\_surveyed + tidal\_height\_m}

\strong{Finding your tidal height:}
\begin{itemize}

\item{} UK: UKHO EasyTide (easytide.ukho.gov.uk) \bsl{}u2014 free 7-day predictions for
any standard port. Download the predicted heights for your survey period.
\item{} Ireland: Marine Institute tide gauges (data.marine.ie)
\item{} France/EU: SHOM (shom.fr) or Copernicus Marine (marine.copernicus.eu)
\item{} Any port: apps such as PocketTides, Tides Near Me, or TidePod give
height-above-CD predictions for named ports.

\end{itemize}


For a typical loch or sheltered bay survey, reading the tidal height from a
tide table for the nearest standard port at the mid-point of the survey is
sufficient. For multi-day surveys or those spanning a full tidal cycle,
supply per-row heights matched to the survey timestamp.
\end{Description}
%
\begin{Usage}
\begin{verbatim}
correct_to_chart_datum(
  df,
  tidal_height_m = NULL,
  depth_col = "depth",
  datetime_col = NULL,
  keep_raw = TRUE,
  verbose = TRUE
)
\end{verbatim}
\end{Usage}
%
\begin{Arguments}
\begin{ldescription}
\item[\code{df}] A dataframe containing at minimum a depth column. Typically the
output of \code{\LinkA{merge\_sensor\_data()}{merge.Rul.sensor.Rul.data}} before passing to \code{\LinkA{predict\_oyster()}{predict.Rul.oyster}}.

\item[\code{tidal\_height\_m}] Numeric scalar or vector. Tidal height above Chart
Datum at the time of survey, in \strong{metres}. If a scalar, the same offset
is applied to all rows. If a vector, must be the same length as \code{nrow(df)}.
Use \code{NULL} to skip correction (a warning will be issued).

\item[\code{depth\_col}] Character. Name of the depth column to correct
(default \code{"depth"}).

\item[\code{datetime\_col}] Character or \code{NULL}. Name of a datetime column in \code{df}.
If supplied, it is used only for informational messages \bsl{}u2014 tidal heights
must still be supplied by the user. Default \code{NULL}.

\item[\code{keep\_raw}] Logical. If \code{TRUE}, retains the original depth values in a
column named \code{depth\_raw\_m} alongside the corrected \code{depth} column
(default \code{TRUE}).

\item[\code{verbose}] Logical. Print correction summary (default \code{TRUE}).
\end{ldescription}
\end{Arguments}
%
\begin{Value}
The input dataframe with the depth column corrected to Chart Datum.
If \code{keep\_raw = TRUE}, the original values are preserved as \code{depth\_raw\_m}.
\end{Value}
%
\begin{Examples}
\begin{ExampleCode}
## Not run: 
# Survey conducted around high water at Oban \u2014 tidal height ~3.1 m above CD
survey_corrected <- correct_to_chart_datum(survey, tidal_height_m = 3.1)

# Per-row heights from a tide gauge CSV matched to survey timestamps
tide_series <- read.csv("oban_tide_gauge.csv")
# (match to survey rows by timestamp \u2014 user responsibility)
survey_corrected <- correct_to_chart_datum(
  survey,
  tidal_height_m = tide_series$height_m
)

## End(Not run)
\end{ExampleCode}
\end{Examples}
\HeaderA{detect\_season}{Detect meteorological season from date and latitude}{detect.Rul.season}
%
\begin{Description}
Returns the meteorological season ("winter", "spring", "summer", "autumn")
for a given date and latitude. Hemisphere is inferred from the sign of
latitude (positive = Northern Hemisphere, negative = Southern Hemisphere).
\end{Description}
%
\begin{Usage}
\begin{verbatim}
detect_season(date, lat)
\end{verbatim}
\end{Usage}
%
\begin{Arguments}
\begin{ldescription}
\item[\code{date}] A \code{Date} or \code{POSIXct} object, or a character string coercible
to \code{Date} (e.g., \code{"2024-03-15"}).

\item[\code{lat}] Numeric. Latitude in decimal degrees. Positive = N hemisphere.
\end{ldescription}
\end{Arguments}
%
\begin{Value}
A character string: one of \code{"winter"}, \code{"spring"}, \code{"summer"},
\code{"autumn"}.
\end{Value}
%
\begin{Examples}
\begin{ExampleCode}
detect_season("2024-01-15", lat = 51.5)   # "winter" (UK)
detect_season("2024-07-15", lat = 51.5)   # "summer"
detect_season("2024-01-15", lat = -33.9)  # "summer" (Sydney)
\end{ExampleCode}
\end{Examples}
\HeaderA{estimate\_chlorophyll\_from\_backscatter}{Estimate chlorophyll-a concentration from ADCP acoustic backscatter}{estimate.Rul.chlorophyll.Rul.from.Rul.backscatter}
%
\begin{Description}
Uses the empirical log-linear relationship between volume backscatter
strength (Sv) and chlorophyll-a concentration to add an estimated
\code{chlorophyll\_a} column to a survey dataframe. Intended for surveys where
a fluorometer was not deployed and the ADCP backscatter intensity is
available.

The conversion is: \AsIs{\texttt{Chl\_a (\bsl{}u00b5g/L) = 10 \textasciicircum{} (a \bsl{}u00b7 Sv + b)}}

Default \code{a} and \code{b} constants are derived from Deines (1999) for three
common ADCP frequencies (300, 600, 1200 kHz). For best results, calibrate
against concurrent water samples from your survey using the \code{calibration}
argument.
\end{Description}
%
\begin{Usage}
\begin{verbatim}
estimate_chlorophyll_from_backscatter(
  df,
  backscatter_col,
  frequency_khz = 600,
  calibration = NULL,
  min_chl = 0.05,
  max_chl = 50,
  keep_raw = TRUE,
  verbose = TRUE
)
\end{verbatim}
\end{Usage}
%
\begin{Arguments}
\begin{ldescription}
\item[\code{df}] Dataframe containing a backscatter column and optional temperature
and depth columns.

\item[\code{backscatter\_col}] Character. Name of the column containing volume
backscatter strength values (Sv in dB re 1 m\bsl{}u207b\bsl{}u00b9). For Nortek Signature
outputs, this is typically the mean amplitude across clean bins, available
as \code{amp\_mean\_dB} after \code{\LinkA{read\_nortek\_adcp()}{read.Rul.nortek.Rul.adcp}}.

\item[\code{frequency\_khz}] Numeric. ADCP operating frequency in kHz. Used to
select default calibration constants if \code{calibration} is not supplied.
Common values: 300, 600, 1200. For other frequencies, supply \code{calibration}
directly.

\item[\code{calibration}] Numeric vector of length 2: \code{c(a, b)} for the log-linear
model \code{log10(Chl\_a) = a * Sv + b}. Supply to override the built-in
frequency-based defaults. Derive from concurrent water samples:
fit \code{log10(observed\_chl) \textasciitilde{} backscatter} and extract intercept (\code{b}) and
slope (\code{a}).

\item[\code{min\_chl}] Numeric. Floor for estimated chlorophyll-a in \bsl{}u00b5g/L
(default \code{0.05}). Prevents physically implausible negative estimates.

\item[\code{max\_chl}] Numeric. Ceiling for estimated chlorophyll-a in \bsl{}u00b5g/L
(default \code{50}). Values above this are likely turbidity artefacts (non-algal
particles dominating backscatter).

\item[\code{keep\_raw}] Logical. If \code{TRUE}, preserves the input backscatter column
(default \code{TRUE}).

\item[\code{verbose}] Logical. Print conversion summary and warnings (default \code{TRUE}).
\end{ldescription}
\end{Arguments}
%
\begin{Value}
The input dataframe with an additional \code{chlorophyll\_a} column
(\bsl{}u00b5g/L). If a \code{chlorophyll\_a} column already exists, a warning is issued
and the existing column is overwritten.
\end{Value}
%
\begin{Section}{Improving accuracy with water samples}

Collect 5\bsl{}u201315 grab/Niskin water samples distributed across your survey area
and depth range, then:

\begin{alltt}# Fit calibration from samples
fit  <- lm(log10(sample_chl) ~ backscatter, data = samples_df)
a    <- coef(fit)[["backscatter"]]
b    <- coef(fit)[["(Intercept)"]]
df   <- estimate_chlorophyll_from_backscatter(
          df, "amp_mean_dB", frequency_khz = 600,
          calibration = c(a, b))
\end{alltt}

\end{Section}
%
\begin{Examples}
\begin{ExampleCode}
## Not run: 
# Default calibration (300 kHz Nortek Signature)
adcp <- read_nortek_adcp("survey.csv")
adcp <- estimate_chlorophyll_from_backscatter(
          adcp, backscatter_col = "amp_mean_dB",
          frequency_khz = 300)

# Custom calibration from water samples
adcp <- estimate_chlorophyll_from_backscatter(
          adcp, "amp_mean_dB", frequency_khz = 300,
          calibration = c(a = 0.041, b = -0.52))

## End(Not run)
\end{ExampleCode}
\end{Examples}
\HeaderA{estimate\_fetch}{Estimate effective fetch from a set of lat/lon points and a bearing}{estimate.Rul.fetch}
%
\begin{Description}
Simple great-circle ray-tracing to estimate how far open water extends
from a point in a given direction, up to \code{max\_fetch\_km}. This is a
lightweight GIS-free approximation \bsl{}u2014 for precise fetch, use QGIS or
the \code{fetchR} R package.
\end{Description}
%
\begin{Usage}
\begin{verbatim}
estimate_fetch(
  lat,
  lon,
  bearing_deg,
  land_polygons = NULL,
  max_fetch_km = 200,
  step_km = 1
)
\end{verbatim}
\end{Usage}
%
\begin{Arguments}
\begin{ldescription}
\item[\code{lat}] Numeric. Site latitude (decimal degrees).

\item[\code{lon}] Numeric. Site longitude (decimal degrees).

\item[\code{bearing\_deg}] Numeric. Prevailing wind direction to check (0 = North).

\item[\code{land\_polygons}] SpatialPolygons or sf POLYGON object representing
land. If NULL, returns max\_fetch\_km (fully open ocean assumed).

\item[\code{max\_fetch\_km}] Numeric. Maximum fetch to search (default 200 km).

\item[\code{step\_km}] Numeric. Ray step size (default 1 km).
\end{ldescription}
\end{Arguments}
%
\begin{Value}
Numeric. Estimated fetch in km.
\end{Value}
\HeaderA{export\_contours}{Export optional contour lines as a GeoPackage for QGIS}{export.Rul.contours}
%
\begin{Description}
Derives contour lines from the IDW suitability raster and saves them as a
vector GeoPackage (\code{.gpkg}). Load alongside the heatmap in QGIS for the
depth-contour look shown in BioBase outputs.
\end{Description}
%
\begin{Usage}
\begin{verbatim}
export_contours(tif_path, interval = 0.1, gpkg_path = NULL)
\end{verbatim}
\end{Usage}
%
\begin{Arguments}
\begin{ldescription}
\item[\code{tif\_path}] Character. Path to the GeoTIFF produced by \code{\LinkA{export\_geotiff()}{export.Rul.geotiff}}.

\item[\code{interval}] Numeric. Contour interval in suitability units (default \code{0.1},
producing lines at 0.1, 0.2, ... 0.9).

\item[\code{gpkg\_path}] Character. Output \code{.gpkg} path. Defaults to same directory
as the \code{.tif} with \AsIs{\texttt{\_contours.gpkg}} suffix.
\end{ldescription}
\end{Arguments}
%
\begin{Value}
The \code{.gpkg} file path (invisibly).
\end{Value}
%
\begin{Examples}
\begin{ExampleCode}
## Not run: 
export_geotiff(result, "oyster_heatmap.tif")
export_contours("oyster_heatmap.tif")

## End(Not run)
\end{ExampleCode}
\end{Examples}
\HeaderA{export\_geotiff}{Export suitability scores as a smooth interpolated GeoTIFF heatmap for QGIS}{export.Rul.geotiff}
%
\begin{Description}
Takes the scored point dataset from \code{\LinkA{predict\_oyster()}{predict.Rul.oyster}} and produces a
\strong{smooth, continuous raster heatmap} using Inverse Distance Weighting (IDW)
interpolation \bsl{}u2014 the same technique used by systems like Lowrance BioBase.
The result looks like a fluid heat surface over your survey area rather than
isolated point squares.

The output GeoTIFF contains three bands:
\begin{enumerate}

\item{} \strong{suitability} \bsl{}u2014 IDW-interpolated score \LinkA{0, 1}{0, 1}. Use this for the heatmap.
\item{} \strong{excluded\_mask} \bsl{}u2014 1 where a survey point was hard-excluded, 0 otherwise.
\item{} \strong{n\_observations} \bsl{}u2014 how many survey points fell within each raster cell
(data density diagnostic).

\end{enumerate}

\end{Description}
%
\begin{Usage}
\begin{verbatim}
export_geotiff(
  df,
  path,
  resolution = 2e-04,
  idw_power = 2,
  idw_max_dist = NULL,
  buffer = 0.001,
  contours = TRUE,
  crs = "EPSG:4326",
  overwrite = TRUE
)
\end{verbatim}
\end{Usage}
%
\begin{Arguments}
\begin{ldescription}
\item[\code{df}] A dataframe processed by \code{\LinkA{predict\_oyster()}{predict.Rul.oyster}}, containing columns
\code{lat}, \code{lon}, \code{suitability}, \code{excluded}.

\item[\code{path}] Character. Output \code{.tif} file path.

\item[\code{resolution}] Numeric. Raster cell size in decimal degrees. Default
\code{0.0002} \bsl{}u2248 22 m \bsl{}u2014 matches the typical 4-decimal-place survey grid. For
large estuary surveys increase to \code{0.001} (100 m); for very detailed
harbour surveys decrease to \code{0.0001} (\textasciitilde{}11 m).

\item[\code{idw\_power}] Numeric. IDW distance decay exponent (default \code{2}). Higher
values make the surface honour nearby points more tightly and decay faster
with distance \bsl{}u2014 producing sharper gradients. Lower values (e.g. 1) give
broader, smoother spread.

\item[\code{idw\_max\_dist}] Numeric or \code{NULL}. Maximum search radius in degrees for
IDW interpolation. If \code{NULL} (default), the radius is automatically set to
5\bsl{}u00d7 the median nearest-neighbour spacing between survey points \bsl{}u2014 so the
heatmap fills gaps between transect lines but does not bleed onto land or
beyond the survey boundary. Override with an explicit value (e.g. \code{0.003}
\bsl{}u2248 330 m) if needed.

\item[\code{buffer}] Numeric. Padding in degrees added around the survey extent
before building the raster grid (default \code{0.001} \bsl{}u2248 110 m). Keeps the
heatmap from clipping at the outermost survey points without extending far
beyond the survey boundary.

\item[\code{contours}] Logical. If \code{TRUE} (default), automatically export contour
lines as a GeoPackage alongside the GeoTIFF using \code{\LinkA{export\_contours()}{export.Rul.contours}}.

\item[\code{crs}] Character. CRS string (default \code{"EPSG:4326"} \bsl{}u2014 WGS84). Use a
projected CRS (e.g. \code{"EPSG:32630"}) if your coordinates are in metres.

\item[\code{overwrite}] Logical. Overwrite existing file (default \code{TRUE}).
\end{ldescription}
\end{Arguments}
%
\begin{Value}
The output file path (invisibly). Also writes a \code{.qml} QGIS style
file alongside the \code{.tif} automatically.
\end{Value}
%
\begin{Examples}
\begin{ExampleCode}
## Not run: 
result <- predict_oyster("survey.csv", "ostrea_edulis")

# Standard export \u2014 smooth heatmap with auto radius and contours
export_geotiff(result, "oyster_heatmap.tif")

# Finer resolution for a small harbour survey
export_geotiff(result, "oyster_heatmap.tif", resolution = 0.0001)

# Override IDW radius explicitly (e.g. sparse transect data)
export_geotiff(result, "oyster_heatmap.tif", idw_max_dist = 0.005)

## End(Not run)
\end{ExampleCode}
\end{Examples}
\HeaderA{export\_qml\_style}{Export a QGIS colour style (.qml) matching the BioBase orange/red heatmap}{export.Rul.qml.Rul.style}
%
\begin{Description}
Writes a QGIS Layer Style file (\code{.qml}) using a yellow \bsl{}u2192 orange \bsl{}u2192 red \bsl{}u2192
dark red colour ramp, closely matching the BioBase/Lowrance heatmap style.
The file is automatically applied when its name matches the \code{.tif} \bsl{}u2014 no
manual styling needed in QGIS.
\end{Description}
%
\begin{Usage}
\begin{verbatim}
export_qml_style(tif_path)
\end{verbatim}
\end{Usage}
%
\begin{Arguments}
\begin{ldescription}
\item[\code{tif\_path}] Character. Path to the \code{.tif} from \code{\LinkA{export\_geotiff()}{export.Rul.geotiff}}.
\end{ldescription}
\end{Arguments}
%
\begin{Value}
The \code{.qml} file path (invisibly).
\end{Value}
\HeaderA{fetch\_live\_environmental\_data}{Fetch live environmental data for a survey extent}{fetch.Rul.live.Rul.environmental.Rul.data}
%
\begin{Description}
Convenience wrapper that fetches data from one or more live sources and
returns a named list of dataframes that can be passed to the relevant
scoring functions. Each source can be fetched independently or all at once.

This function requires internet access and registered API credentials where
applicable (see \code{\LinkA{oystermapR\_live\_config()}{oystermapR.Rul.live.Rul.config}}). If a source fails, it is
silently skipped and not included in the output \bsl{}u2014 the function never aborts
due to a single failed source.
\end{Description}
%
\begin{Usage}
\begin{verbatim}
fetch_live_environmental_data(
  survey,
  sources = "all",
  date_range = NULL,
  verbose = TRUE
)
\end{verbatim}
\end{Usage}
%
\begin{Arguments}
\begin{ldescription}
\item[\code{survey}] Dataframe with \code{lat} and \code{lon} columns defining the survey
extent. The bounding box of these coordinates is used as the query extent.

\item[\code{sources}] Character vector. Which data sources to fetch. Options:
\code{"cmems"}, \code{"ices\_hab"}, \code{"ices\_vms"}, \code{"emodnet\_predators"},
\code{"fsa\_shellfish"}. Use \code{"all"} to attempt all sources.

\item[\code{date\_range}] Character vector of length 2: \code{c("YYYY-MM-DD", "YYYY-MM-DD")}.
Date range for time-windowed queries (CMEMS, ICES HAB). Defaults to
the last 5 years.

\item[\code{verbose}] Logical. Print progress (default TRUE).
\end{ldescription}
\end{Arguments}
%
\begin{Value}
Named list. Each successfully fetched source produces one element:
\begin{itemize}

\item{} \AsIs{\texttt{\$cmems}}: SST, salinity, chlorophyll grid
\item{} \AsIs{\texttt{\$ices\_hab}}: HAB event records within extent
\item{} \AsIs{\texttt{\$ices\_vms}}: Swept-area ratio per ICES c-square
\item{} \AsIs{\texttt{\$emodnet\_predators}}: Predator occurrence records
\item{} \AsIs{\texttt{\$fsa\_shellfish}}: Shellfish classification polygons / area descriptions

\end{itemize}

\end{Value}
%
\begin{Examples}
\begin{ExampleCode}
## Not run: 
oystermapR_live_config(cmems_user = "u", cmems_password = "p")
live <- fetch_live_environmental_data(survey, sources = c("ices_hab", "ices_vms"))
result <- score_hab_risk(result, hab_data = live$ices_hab, species = "ostrea_edulis")
result <- score_anthropogenic_disturbance(result, trawling_data = live$ices_vms)

## End(Not run)
\end{ExampleCode}
\end{Examples}
\HeaderA{generate\_report}{Generate a formatted survey report from predict\_oyster() results}{generate.Rul.report}
%
\begin{Description}
Renders a structured assessment report from the output of \code{\LinkA{predict\_oyster()}{predict.Rul.oyster}}.
The report includes an executive summary, suitability class breakdown,
per-variable scoring table, most common limiting factors, top introduction
sites, and a methodology section.

Output format is PDF by default (requires a LaTeX installation \bsl{}u2014 see Note
below). HTML is also supported and requires no additional system software.
\end{Description}
%
\begin{Usage}
\begin{verbatim}
generate_report(
  result,
  output = "oyster_report.html",
  title = "Oyster Suitability Assessment",
  author = "",
  species = "ostrea_edulis",
  top_n = 5L,
  validation = NULL,
  comparison = NULL,
  open = TRUE
)
\end{verbatim}
\end{Usage}
%
\begin{Arguments}
\begin{ldescription}
\item[\code{result}] A dataframe returned by \code{\LinkA{predict\_oyster()}{predict.Rul.oyster}}.

\item[\code{output}] Character. Output file path. Extension determines format:
\code{.pdf} for PDF, \code{.html} for HTML. Default \code{"oyster\_report.pdf"}.

\item[\code{title}] Character. Report title. Default \code{"Oyster Suitability Assessment"}.

\item[\code{author}] Character. Author name(s) to appear on the title page.
Default \code{""} (omitted).

\item[\code{species}] Character. Species key used in \code{\LinkA{predict\_oyster()}{predict.Rul.oyster}}.
Default \code{"ostrea\_edulis"}.

\item[\code{top\_n}] Integer. Number of top introduction sites to include in the
report (default \code{5}).

\item[\code{validation}] Named list or \code{NULL}. Optional output from
\code{\LinkA{validate\_against\_records()}{validate.Rul.against.Rul.records}}. When supplied, adds a Validation section
with ROC curve, AUC, TSS, and confusion-matrix metrics.

\item[\code{comparison}] Dataframe or \code{NULL}. Optional output from
\code{\LinkA{compare\_species()}{compare.Rul.species}}. When supplied, adds a Competition Adjustment section
showing the \emph{M. gigas} competitive pressure penalty on \emph{O. edulis} scores.

\item[\code{open}] Logical. Open the report in your default viewer after rendering
(default \code{TRUE}).
\end{ldescription}
\end{Arguments}
%
\begin{Value}
The output file path (invisibly).
\end{Value}
%
\begin{Note}
\strong{HTML reports} (recommended) require \code{plotly} and \code{scales}:
\code{install.packages(c("plotly","scales"))}.
\strong{PDF reports} additionally require a LaTeX installation; see the
Note section below.

\strong{PDF output requires LaTeX.} If you don't have LaTeX installed,
either use \code{output = "report.html"} (no extra software needed), or install
a minimal LaTeX distribution with:

\begin{alltt}install.packages("tinytex")
tinytex::install_tinytex()
\end{alltt}

\end{Note}
%
\begin{Examples}
\begin{ExampleCode}
## Not run: 
result <- predict_oyster(survey, "ostrea_edulis",
                         output_geotiff = "kames_bay.tif")

# Standard HTML report (recommended \u2014 Plotly charts, no LaTeX needed)
generate_report(result, "kames_bay_report.html",
                title  = "Kames Bay Oyster Habitat Assessment",
                author = "T. Tucker")

# With validation results embedded
val  <- validate_against_records(result, nbn_records)
generate_report(result, "kames_bay_validated.html",
                validation = val)

# With competition adjustment (when comparing species)
comp <- compare_species(survey,
          species = c("ostrea_edulis", "magallana_gigas"))
generate_report(result, "kames_bay_competition.html",
                comparison = comp)

# PDF report (requires LaTeX / tinytex)
generate_report(result, "kames_bay_report.pdf",
                title = "Kames Bay Oyster Habitat Assessment")

## End(Not run)
\end{ExampleCode}
\end{Examples}
\HeaderA{get\_species\_tolerances}{Get species tolerance parameters}{get.Rul.species.Rul.tolerances}
%
\begin{Description}
Get species tolerance parameters
\end{Description}
%
\begin{Usage}
\begin{verbatim}
get_species_tolerances(species)
\end{verbatim}
\end{Usage}
%
\begin{Arguments}
\begin{ldescription}
\item[\code{species}] Character. Species key (e.g. \code{"ostrea\_edulis"}), latin name,
or common name (partial, case-insensitive).
\end{ldescription}
\end{Arguments}
%
\begin{Value}
Named list of tolerance parameters for the species.
\end{Value}
%
\begin{Examples}
\begin{ExampleCode}
tol <- get_species_tolerances("ostrea_edulis")
tol$exclusions$temperature
\end{ExampleCode}
\end{Examples}
\HeaderA{get\_tolerance\_posteriors}{Retrieve current posterior tolerance parameters for a species}{get.Rul.tolerance.Rul.posteriors}
%
\begin{Description}
Returns the tolerance parameter list currently held in the session cache,
after any \code{\LinkA{update\_species\_tolerances()}{update.Rul.species.Rul.tolerances}} calls. If no update has been run,
returns the built-in prior parameters unchanged.
\end{Description}
%
\begin{Usage}
\begin{verbatim}
get_tolerance_posteriors(species, var = NULL)
\end{verbatim}
\end{Usage}
%
\begin{Arguments}
\begin{ldescription}
\item[\code{species}] Character. Species key (e.g. \code{"ostrea\_edulis"}).

\item[\code{var}] Character or NULL. If specified, returns parameters for that
variable only; otherwise returns the full \AsIs{\texttt{\$scored}} list.
\end{ldescription}
\end{Arguments}
%
\begin{Value}
Named list of tolerance parameters (same structure as
\code{\LinkA{get\_species\_tolerances()}{get.Rul.species.Rul.tolerances}}\AsIs{\texttt{\$scored}}).
\end{Value}
%
\begin{Examples}
\begin{ExampleCode}
## Not run: 
# After calling update_species_tolerances()
post <- get_tolerance_posteriors("ostrea_edulis")
post$temperature  # updated optimal_min, optimal_max with CIs

## End(Not run)
\end{ExampleCode}
\end{Examples}
\HeaderA{identify\_resilient\_sites}{Identify climate-resilient locations}{identify.Rul.resilient.Rul.sites}
%
\begin{Description}
From a \code{\LinkA{project\_suitability()}{project.Rul.suitability}} output, identifies locations that maintain
at least \code{min\_class} suitability across all projected scenarios. These are
the most robust sites for long-term restoration or aquaculture investment.
\end{Description}
%
\begin{Usage}
\begin{verbatim}
identify_resilient_sites(projections, baseline, min_class = "Moderate")
\end{verbatim}
\end{Usage}
%
\begin{Arguments}
\begin{ldescription}
\item[\code{projections}] Named list returned by \code{\LinkA{project\_suitability()}{project.Rul.suitability}}.

\item[\code{baseline}] Dataframe from \code{\LinkA{predict\_oyster()}{predict.Rul.oyster}} (the baseline result).

\item[\code{min\_class}] Character. Minimum suitability class required in all
scenarios. One of \code{"High"}, \code{"Moderate"}, \code{"Low"} (default \code{"Moderate"}).
\end{ldescription}
\end{Arguments}
%
\begin{Value}
Subset of \code{baseline} for locations that meet \code{min\_class} in every
scenario, with an added \code{n\_scenarios\_qualifying} column.
\end{Value}
%
\begin{Examples}
\begin{ExampleCode}
## Not run: 
proj      <- project_suitability(result, tol)
resilient <- identify_resilient_sites(proj, result, min_class = "Moderate")
nrow(resilient)  # sites that stay Moderate+ even under RCP8.5

## End(Not run)
\end{ExampleCode}
\end{Examples}
\HeaderA{interpolate\_survey}{Interpolate survey measurements to a regular grid before scoring}{interpolate.Rul.survey}
%
\begin{Description}
Converts sparse point observations (CTD casts, ADCP profiles) to a regular
grid by spatial interpolation. The interpolated grid is returned in the same
dataframe format as the raw survey, and can be passed directly to
\code{\LinkA{predict\_oyster()}{predict.Rul.oyster}}.

\strong{Method selection} (applied per variable independently):
\begin{itemize}

\item{} \strong{Ordinary kriging} is used when \bsl{}u2265 50 non-NA observations are available
for a variable, the \code{gstat} and \code{sp} packages are installed, and the
variogram can be fitted without error. Kriging is the best linear
unbiased estimator (BLUE) under the assumption of spatial stationarity \bsl{}u2014
it minimises mean squared prediction error and returns a kriging variance
for each cell.
\item{} \strong{IDW} (Inverse Distance Weighting, p = 2) is used as a fallback when
kriging cannot be applied. It is fast, deterministic, and requires no
model fitting, but produces no uncertainty estimate and can create
bull's-eye artefacts around isolated observations.

\end{itemize}


Kriging variograms are fitted automatically using \code{gstat::fit.variogram()}
with a Mat\bsl{}u00e9rn model (default) \bsl{}u2014 the most flexible and statistically
justified model for environmental data. If automatic fitting fails, a
spherical model is tried; if that also fails, IDW is used.
\end{Description}
%
\begin{Usage}
\begin{verbatim}
interpolate_survey(
  survey,
  vars = NULL,
  resolution_m = 100,
  method = c("auto", "kriging", "idw"),
  kriging_min_n = 50L,
  idw_power = 2,
  idw_max_radius_m = NULL,
  verbose = TRUE
)
\end{verbatim}
\end{Usage}
%
\begin{Arguments}
\begin{ldescription}
\item[\code{survey}] Dataframe with \code{lat}, \code{lon}, and measurement columns.

\item[\code{vars}] Character vector of column names to interpolate. Default: all
numeric columns except \code{lat} and \code{lon}.

\item[\code{resolution\_m}] Numeric. Grid cell size in metres (default 100 m).
Coarsen for sparse surveys (> 500 m) or refine for dense ADCP grids (50 m).

\item[\code{method}] Character. \code{"auto"} (default \bsl{}u2014 kriging if possible, IDW
fallback), \code{"kriging"} (force kriging; error if insufficient data), or
\code{"idw"} (always use IDW).

\item[\code{kriging\_min\_n}] Integer. Minimum observations required to attempt
kriging (default 50).

\item[\code{idw\_power}] Numeric. IDW distance decay exponent (default 2).

\item[\code{idw\_max\_radius\_m}] Numeric. Maximum search radius for IDW neighbours
in metres (default 5 * resolution\_m).

\item[\code{verbose}] Logical. Print per-variable method used and diagnostics
(default TRUE).
\end{ldescription}
\end{Arguments}
%
\begin{Value}
A dataframe on a regular grid with interpolated values, a \code{method}
column per variable (\AsIs{\texttt{interp\_method\_<var>}}), and kriging variance columns
(\AsIs{\texttt{krige\_var\_<var>}}) where kriging was used.
\end{Value}
%
\begin{Examples}
\begin{ExampleCode}
## Not run: 
# Auto method \u2014 kriging for dense variables, IDW for sparse
grid <- interpolate_survey(survey, resolution_m = 100)
result <- predict_oyster(grid, "ostrea_edulis")

# Force IDW (fast, no gstat needed)
grid <- interpolate_survey(survey, resolution_m = 200, method = "idw")

# Inspect which method was used per variable
unique(grid$interp_method_temperature)

# Kriging variance \u2014 high values = uncertain interpolation
hist(grid$krige_var_temperature)

## End(Not run)
\end{ExampleCode}
\end{Examples}
\HeaderA{list\_species}{List all supported species}{list.Rul.species}
%
\begin{Description}
List all supported species
\end{Description}
%
\begin{Usage}
\begin{verbatim}
list_species()
\end{verbatim}
\end{Usage}
%
\begin{Value}
Named character vector: names = latin names, values = species keys.
\end{Value}
%
\begin{Examples}
\begin{ExampleCode}
list_species()
\end{ExampleCode}
\end{Examples}
\HeaderA{load\_tolerance\_update}{Load saved Bayesian tolerance updates into session cache}{load.Rul.tolerance.Rul.update}
%
\begin{Description}
Reads a previously saved \code{.rds} file (from \code{\LinkA{save\_tolerance\_update()}{save.Rul.tolerance.Rul.update}}) and
restores it into the session cache. Subsequent calls to \code{\LinkA{predict\_oyster()}{predict.Rul.oyster}}
and \code{\LinkA{score\_locations()}{score.Rul.locations}} will automatically use the loaded parameters.
\end{Description}
%
\begin{Usage}
\begin{verbatim}
load_tolerance_update(species = "all", path = "~/.oystermapR", verbose = TRUE)
\end{verbatim}
\end{Usage}
%
\begin{Arguments}
\begin{ldescription}
\item[\code{species}] Character. Species key or \code{"all"} to load all \code{.rds} files
in \code{path}.

\item[\code{path}] Character. Directory to load from (default: \AsIs{\texttt{\textasciitilde{}/.oystermapR/}}).

\item[\code{verbose}] Logical. Default TRUE.
\end{ldescription}
\end{Arguments}
%
\begin{Value}
Invisibly: the loaded tolerance list.
\end{Value}
%
\begin{Examples}
\begin{ExampleCode}
## Not run: 
load_tolerance_update("ostrea_edulis")
# Now predict_oyster() uses the updated tolerances automatically
result <- predict_oyster(survey, "ostrea_edulis")

## End(Not run)
\end{ExampleCode}
\end{Examples}
\HeaderA{merge\_sensor\_data}{Merge multiple sensor datasets into a single survey table}{merge.Rul.sensor.Rul.data}
%
\begin{Description}
Takes any number of sensor dataframes (from \code{\LinkA{read\_nortek\_adcp()}{read.Rul.nortek.Rul.adcp}},
\code{\LinkA{read\_generic\_csv()}{read.Rul.generic.Rul.csv}}, or custom sources) and merges them spatially into a
single table ready for \code{\LinkA{predict\_oyster()}{predict.Rul.oyster}}.

Each dataset is first rounded to a common spatial resolution grid. Datasets
are then joined using the grid cell key \bsl{}u2014 cells that exist in multiple
datasets are merged; cells that exist in only one dataset carry \code{NA} for
variables from other sensors.
\end{Description}
%
\begin{Usage}
\begin{verbatim}
merge_sensor_data(..., spatial_res = 4L, date_source = NULL, verbose = TRUE)
\end{verbatim}
\end{Usage}
%
\begin{Arguments}
\begin{ldescription}
\item[\code{...}] Two or more dataframes to merge. Each must contain \code{lat} and \code{lon}
columns. Pass them as named arguments for cleaner messages (e.g.
\AsIs{\texttt{adcp = adcp\_df, biobase = biobase\_df}}).

\strong{Multiple runs of the same sensor:} If you have several surveys from the
same sensor type, pass them as a list and they will be automatically
stacked via \code{\LinkA{stack\_surveys()}{stack.Rul.surveys}} before merging:

\begin{alltt}merge_sensor_data(
  adcp    = list(adcp_run1, adcp_run2, adcp_run3),
  bathy   = bathy_df,
  biobase = biobase_df
)
\end{alltt}


\item[\code{spatial\_res}] Integer. Decimal places for the common spatial grid
(default \code{4}, \bsl{}u224811 m at 56\bsl{}u00b0N). Must be the same or coarser than the
resolution used when reading each individual sensor.

\item[\code{date\_source}] Character. Name of the dataset to take the \code{date} column
from (default \code{NULL} \bsl{}u2014 uses the first dataset that contains a \code{date}
column).

\item[\code{verbose}] Logical. Print a merge summary (default \code{TRUE}).
\end{ldescription}
\end{Arguments}
%
\begin{Value}
A single dataframe with all oystermapR-compatible columns populated
from their respective sensor sources. Columns that weren't available from
any sensor will be absent from the result \bsl{}u2014 \code{\LinkA{predict\_oyster()}{predict.Rul.oyster}} handles
missing columns gracefully.
\end{Value}
%
\begin{Examples}
\begin{ExampleCode}
## Not run: 
# Single run per sensor
adcp    <- read_nortek_adcp("adcp.csv")
biobase <- read_generic_csv("biobase.csv")
survey  <- merge_sensor_data(adcp = adcp, biobase = biobase)

# Multiple ADCP runs (different survey days) \u2014 automatically stacked
adcp1  <- read_nortek_adcp("survey_jan.csv")
adcp2  <- read_nortek_adcp("survey_feb.csv")
survey <- merge_sensor_data(adcp = list(adcp1, adcp2), biobase = biobase)

result <- predict_oyster(survey, "ostrea_edulis", output_geotiff = "map.tif")

## End(Not run)
\end{ExampleCode}
\end{Examples}
\HeaderA{oystermapR\_help}{Print an interactive getting-started guide for oystermapR}{oystermapR.Rul.help}
%
\begin{Description}
Prints a step-by-step workflow guide to the console, covering sensor data
ingestion, spatial merging, suitability prediction, and QGIS export.
Run this any time you need a reminder of the available functions and
typical usage patterns.
\end{Description}
%
\begin{Usage}
\begin{verbatim}
oystermapR_help(topic = NULL)
\end{verbatim}
\end{Usage}
%
\begin{Arguments}
\begin{ldescription}
\item[\code{topic}] Character. Optionally focus on a specific topic:
\code{"sensors"}, \code{"merge"}, \code{"predict"}, \code{"qgis"}, \code{"species"}.
Default \code{NULL} prints the full guide.
\end{ldescription}
\end{Arguments}
%
\begin{Value}
Called for its side-effect (console output). Returns \code{invisible(NULL)}.
\end{Value}
%
\begin{Examples}
\begin{ExampleCode}
oystermapR_help()                  # full guide
oystermapR_help("sensors")         # sensor ingestion only
oystermapR_help("qgis")            # QGIS output only
\end{ExampleCode}
\end{Examples}
\HeaderA{oystermapR\_live\_config}{Configure live data API credentials for oystermapR}{oystermapR.Rul.live.Rul.config}
%
\begin{Description}
Stores API credentials in R session options so live data fetching can
authenticate with external services. Credentials are \strong{not} written to
disk and are lost when the R session ends.

All parameters are optional \bsl{}u2014 only set credentials for the services you
intend to use. Functions that require a missing credential will warn
and fall back to manual data rather than aborting.
\end{Description}
%
\begin{Usage}
\begin{verbatim}
oystermapR_live_config(
  cmems_user = NULL,
  cmems_password = NULL,
  ices_key = NULL,
  show_config = FALSE
)
\end{verbatim}
\end{Usage}
%
\begin{Arguments}
\begin{ldescription}
\item[\code{cmems\_user}] Character. CMEMS username.

\item[\code{cmems\_password}] Character. CMEMS password.

\item[\code{ices\_key}] Character. ICES API key (currently optional; ICES APIs are
largely open).

\item[\code{show\_config}] Logical. Print current configuration (credentials masked).
Default FALSE.
\end{ldescription}
\end{Arguments}
%
\begin{Value}
Invisibly returns a named list of the options set.
\end{Value}
%
\begin{Section}{Free registrations}

\begin{itemize}

\item{} \strong{CMEMS}: Register at \Rhref{https://marine.copernicus.eu}{marine.copernicus.eu}
\item{} \strong{ICES}: API access is open (no key required for HAB/VMS endpoints)
\item{} \strong{EMODnet Biology}: Open WFS, no key required
\item{} \strong{FSA/DAERA}: Open REST API, no key required

\end{itemize}

\end{Section}
%
\begin{Examples}
\begin{ExampleCode}
## Not run: 
oystermapR_live_config(
  cmems_user     = "myusername",
  cmems_password = "mypassword"
)
# Then use fetch_live = TRUE in any supporting function
score_hab_risk(result, species = "ostrea_edulis", fetch_live = TRUE)

## End(Not run)
\end{ExampleCode}
\end{Examples}
\HeaderA{parse\_opendrift\_connectivity}{Parse an OpenDrift particle tracking output into a connectivity matrix}{parse.Rul.opendrift.Rul.connectivity}
%
\begin{Usage}
\begin{verbatim}
parse_opendrift_connectivity(
  opendrift_csv,
  bin_deg = 0.05,
  status_col = NULL,
  status_keep = c("stranded", "settled"),
  min_weight = 0.005,
  verbose = TRUE
)
\end{verbatim}
\end{Usage}
%
\begin{Arguments}
\begin{ldescription}
\item[\code{opendrift\_csv}] Character. Path to OpenDrift settlement CSV file. Must
contain columns \code{origin\_lat}, \code{origin\_lon}, \code{final\_lat}, \code{final\_lon}.

\item[\code{bin\_deg}] Numeric. Grid cell size in decimal degrees for binning
(default 0.05 \bsl{}u2248 5 km at mid-latitudes).

\item[\code{status\_col}] Character or NULL. Name of a status column to filter on.
If NULL (default), all rows are used.

\item[\code{status\_keep}] Character. Value(s) of \code{status\_col} to include
(default \code{c("stranded","settled")}).

\item[\code{min\_weight}] Numeric. Minimum connectivity weight to include in output
(default 0.005 = 0.5\bsl{}

\bsl{}itemverboseLogical. Default TRUE.

\end{ldescription}

A dataframe with columns \code{source\_lat}, \code{source\_lon}, \code{dest\_lat},
\code{dest\_lon}, \code{weight} \bsl{}u2014 ready for passing to
\code{\LinkA{score\_larval\_connectivity()}{score.Rul.larval.Rul.connectivity}}.


Converts the output of an OpenDrift settlement simulation (a CSV file with
particle origin and final position) into the connectivity matrix format
expected by \code{score\_larval\_connectivity(connectivity\_matrix = ...)}.

\strong{OpenDrift setup:}
Run OpenDrift (\code{OceanDrift} or \code{OpenOil} adapted for larvae) seeding particles
at each source reef location. Export a CSV with at minimum:
\begin{itemize}

\item{} \code{origin\_lat}, \code{origin\_lon} \bsl{}u2014 seed position
\item{} \code{final\_lat}, \code{final\_lon} \bsl{}u2014 settlement position (stranded or settled)
\item{} \code{status} \bsl{}u2014 filter to \code{"stranded"} or \code{"active"} as appropriate

\end{itemize}


The function bins particles into a grid and computes the fraction of particles
from each source bin that arrive in each destination bin.


## Not run: 
\# After running OpenDrift simulation and exporting CSV:
cm <- parse\_opendrift\_connectivity("opendrift\_output.csv", bin\_deg = 0.05)

result <- predict\_oyster(survey, "ostrea\_edulis")
result <- score\_larval\_connectivity(result,
            species             = "ostrea\_edulis",
            connectivity\_matrix = cm)

## End(Not run)

\end{Arguments}
\HeaderA{permutation\_importance}{Permutation variable importance for suitability model}{permutation.Rul.importance}
%
\begin{Description}
Estimates the importance of each scored environmental variable by randomly
permuting its values and measuring the resulting drop in AUC against known
presence/absence records. A large drop indicates the model strongly depends
on that variable; a small drop indicates it could be removed with minimal
impact.

This approach (Breiman 2001) is model-agnostic, requires no model refit,
and is directly interpretable: "how much does AUC fall if I destroy the
information in variable X?"
\end{Description}
%
\begin{Usage}
\begin{verbatim}
permutation_importance(
  predicted,
  records,
  presence_col = "presence",
  n_permutations = 50L,
  match_radius_deg = 0.002,
  seed = 42L,
  verbose = TRUE
)
\end{verbatim}
\end{Usage}
%
\begin{Arguments}
\begin{ldescription}
\item[\code{predicted}] Dataframe from \code{\LinkA{predict\_oyster()}{predict.Rul.oyster}} with \code{lat}, \code{lon},
\code{suitability}, and per-variable score columns (e.g. \code{score\_temperature},
\code{score\_salinity}). The function uses the per-variable component scores
already computed by \code{\LinkA{predict\_oyster()}{predict.Rul.oyster}}.

\item[\code{records}] Dataframe of presence/absence records (same format as
\code{\LinkA{validate\_against\_records()}{validate.Rul.against.Rul.records}}).

\item[\code{presence\_col}] Character. Name of presence/absence column (default
\code{"presence"}).

\item[\code{n\_permutations}] Integer. Number of permutation replicates per variable
(default 50). More = more stable importance estimates; \bsl{}u2265 100 recommended
for publication.

\item[\code{match\_radius\_deg}] Numeric. Matching radius (default 0.002\bsl{}u00b0).

\item[\code{seed}] Integer. Random seed (default 42).

\item[\code{verbose}] Logical. Default TRUE.
\end{ldescription}
\end{Arguments}
%
\begin{Value}
A dataframe (sorted by importance, descending) with columns:
\begin{itemize}

\item{} \code{variable}: scored variable name
\item{} \code{baseline\_auc}: AUC of unperturbed model
\item{} \code{mean\_permuted\_auc}: mean AUC after permuting this variable
\item{} \code{importance}: baseline\_auc \bsl{}u2212 mean\_permuted\_auc (drop in AUC)
\item{} \code{importance\_sd}: SD of AUC drop across permutation replicates
\item{} \code{importance\_pct}: importance as percentage of baseline AUC
\item{} \code{rank}: rank by importance (1 = most important)

\end{itemize}

\end{Value}
%
\begin{Section}{Interpretation}

\begin{itemize}

\item{} importance > 0.10 (> 10\% AUC drop): high importance \bsl{}u2014 variable is load-
bearing for model discrimination.
\item{} importance 0.02\bsl{}u20130.10: moderate importance.
\item{} importance < 0.02: low importance \bsl{}u2014 consider whether this variable adds
value relative to its data collection cost.

\end{itemize}

\end{Section}
%
\begin{References}
Breiman (2001) Machine Learning 45:5-32.
Zurell et al. (2020) Ecography 43:1261-1277.
\end{References}
%
\begin{Examples}
\begin{ExampleCode}
## Not run: 
records <- read.csv("nbn_ostrea_edulis.csv")
result  <- predict_oyster(survey, "ostrea_edulis")

imp <- permutation_importance(result, records, n_permutations = 100)
print(imp)
# Variable        Importance  Rank
# temperature     0.183       1
# salinity        0.091       2
# depth           0.045       3
# ...

## End(Not run)
\end{ExampleCode}
\end{Examples}
\HeaderA{predict\_oyster}{Predict oyster growth suitability from environmental survey data}{predict.Rul.oyster}
%
\begin{Description}
The primary user-facing function of \code{oystermapR}. Accepts a tabular dataset
(CSV file path or dataframe) containing spatial coordinates and environmental
measurements, applies species-specific exclusion and weighted scoring rules,
and returns a scored dataframe alongside an optional GeoTIFF heatmap for
QGIS visualisation.
\end{Description}
%
\begin{Usage}
\begin{verbatim}
predict_oyster(
  data,
  species,
  output_geotiff = FALSE,
  resolution = 2e-04,
  contours = TRUE,
  verbose = FALSE
)
\end{verbatim}
\end{Usage}
%
\begin{Arguments}
\begin{ldescription}
\item[\code{data}] A dataframe or a file path to a CSV file. Must contain at minimum:
\begin{itemize}

\item{} \strong{\code{lon}} / \strong{\code{lng}} / \strong{\code{longitude}} \bsl{}u2014 longitude in decimal degrees
\item{} \strong{\code{lat}} / \strong{\code{latitude}} \bsl{}u2014 latitude in decimal degrees
\item{} \strong{\code{date}} \bsl{}u2014 observation date (any format coercible by \code{as.Date()})
Additional environmental columns are matched automatically; see
\strong{Column Naming} below.

\end{itemize}


\item[\code{species}] Character string identifying the target oyster species.
Accepts the key (\code{"ostrea\_edulis"}), latin name, or common name.
Use \code{\LinkA{list\_species()}{list.Rul.species}} to see all available options.

\item[\code{output\_geotiff}] Logical or character. If \code{TRUE}, exports a GeoTIFF to
the current working directory as \AsIs{\texttt{<species>\_suitability.tif}}. If a
character string, it is used as the file path. If \code{FALSE} (default), no
raster file is written.

\item[\code{resolution}] Numeric. Raster cell size in degrees for GeoTIFF output
(default \code{0.001}, approximately 100 m at mid-latitudes). Ignored when
\code{output\_geotiff = FALSE}.

\item[\code{verbose}] Logical. If \code{TRUE}, prints a per-variable scoring summary
after processing (default \code{FALSE}).
\end{ldescription}
\end{Arguments}
%
\begin{Value}
A dataframe (invisibly) with all original columns plus:
\begin{itemize}

\item{} \strong{\code{season}}: detected season at each location/date.
\item{} \strong{\code{excluded}}: logical, \code{TRUE} if the location fails a hard exclusion.
\item{} \strong{\code{exclusion\_reason}}: character, reason(s) for exclusion or \code{NA}.
\item{} \strong{\code{suitability}}: numeric \LinkA{0, 1}{0, 1}, overall weighted suitability score.
\item{} \strong{\code{suitability\_class}}: factor, one of "High", "Moderate", "Low",
"Very Low", "Excluded".
\item{} \strong{\AsIs{\texttt{score\_<variable>}}}: per-variable component scores.

\end{itemize}


If \code{output\_geotiff} is set, a GeoTIFF is written to disk and its path is
printed to the console.
\end{Value}
%
\begin{Section}{Column Naming}

Column names are matched case-insensitively. Recognised synonyms:
\Tabular{ll}{
Variable & Recognised column names \\{}
Longitude & \code{lon}, \code{lng}, \code{longitude}, \code{x} \\{}
Latitude & \code{lat}, \code{latitude}, \code{y} \\{}
Date & \code{date}, \code{datetime}, \code{timestamp} \\{}
Temperature (\bsl{}u00b0C) & \code{temperature}, \code{temp}, \code{temp\_c} \\{}
Salinity (PSU) & \code{salinity}, \code{sal}, \code{salinity\_psu} \\{}
Dissolved oxygen (mg/L) & \code{dissolved\_oxygen}, \code{do}, \code{do\_mgl}, \code{oxygen} \\{}
Depth (m) & \code{depth}, \code{depth\_m} \\{}
Current velocity (m/s) & \code{current\_velocity}, \code{velocity}, \code{current}, \code{u\_mean} \\{}
Shear stress (N/m\bsl{}u00b2) & \code{shear\_stress}, \code{tau}, \code{bed\_shear}, \code{shear} \\{}
Chlorophyll-a (\bsl{}u00b5g/L) & \code{chlorophyll\_a}, \code{chla}, \code{chl\_a}, \code{chlorophyll} \\{}
Turbidity (NTU) & \code{turbidity}, \code{ntu}, \code{turb} \\{}
Slope (degrees) & \code{slope}, \code{slope\_deg} \\{}
Roughness / rugosity & \code{roughness}, \code{rugosity} \\{}
Substrate hardness & \code{substrate\_hardness}, \code{hardness}, \code{bottom\_hardness} \\{}
Sediment type & \code{sediment\_type}, \code{sediment}, \code{substrate\_type} \\{}
Benthic communities & \code{benthic\_communities}, \code{benthic}, \code{community} \\{}
Biotope & \code{biotope}, \code{biotopes}, \code{habitat} \\{}
Fishing intensity & \code{fishing\_intensity}, \code{fishing}, \code{fishing\_observed} \\{}
}
\end{Section}
%
\begin{Examples}
\begin{ExampleCode}
## Not run: 
# Using a CSV file
result <- predict_oyster(
  data    = "my_survey.csv",
  species = "ostrea_edulis",
  output_geotiff = "oyster_suitability.tif"
)

# Using a dataframe
df <- read.csv("survey_data.csv")
result <- predict_oyster(df, species = "ostrea_edulis", verbose = TRUE)

# View top locations
subset(result, suitability_class == "High")

## End(Not run)
\end{ExampleCode}
\end{Examples}
\HeaderA{project\_suitability}{Project suitability under future climate scenarios}{project.Rul.suitability}
%
\begin{Description}
Re-scores survey locations under user-specified climate perturbations to
assess how habitat suitability may change under warming and related
oceanographic shifts. Perturbations are applied as additive deltas to
temperature and salinity, and multiplicative factors to other variables.

Built-in scenarios align with UKCP18 marine projections for NW European
shelf seas:
\Tabular{llll}{
Scenario & SST delta & Salinity delta & Notes \\{}
RCP2.6/SSP1 & +0.5\bsl{}u00b0C & 0.0 PSU & Low emissions, best case \\{}
RCP4.5/SSP2 & +1.2\bsl{}u00b0C & -0.2 PSU & Intermediate emissions \\{}
RCP8.5/SSP5 & +2.5\bsl{}u00b0C & -0.5 PSU & High emissions, worst case \\{}
Custom & user & user & Supply your own deltas \\{}
}


The function runs \code{score\_locations()} for each scenario and returns a list
with one result dataframe per scenario, plus a comparison summary and a
\code{suitability\_delta} column (projected minus baseline) for each.
\end{Description}
%
\begin{Usage}
\begin{verbatim}
project_suitability(
  result,
  tolerances,
  scenarios = c("rcp26", "rcp45", "rcp85"),
  custom_deltas = NULL,
  horizon = "2050s",
  verbose = TRUE
)
\end{verbatim}
\end{Usage}
%
\begin{Arguments}
\begin{ldescription}
\item[\code{result}] Dataframe from \code{\LinkA{predict\_oyster()}{predict.Rul.oyster}} (the baseline).

\item[\code{tolerances}] Species tolerance list from \code{\LinkA{get\_species\_tolerances()}{get.Rul.species.Rul.tolerances}}.

\item[\code{scenarios}] Character vector of built-in scenario names, or NULL to
use \code{custom\_deltas} only. Built-in: \code{"rcp26"}, \code{"rcp45"}, \code{"rcp85"}.
Default: all three.

\item[\code{custom\_deltas}] Named list of custom scenarios. Each element is a
named numeric vector of deltas keyed to column names, e.g.
\code{list(my\_scenario = c(temperature = 1.5, salinity = -0.3))}.

\item[\code{horizon}] Character. Label for the time horizon, used in output names
only (default \code{"2050s"}).

\item[\code{verbose}] Logical. Print scenario comparison table (default TRUE).
\end{ldescription}
\end{Arguments}
%
\begin{Value}
Named list:
\begin{itemize}

\item{} One dataframe per scenario (named by scenario key) with \code{suitability},
\code{suitability\_class}, and \code{suitability\_delta} columns.
\item{} \code{"summary"} \bsl{}u2014 dataframe comparing mean suitability, class breakdown,
and net change vs baseline across all scenarios.

\end{itemize}

\end{Value}
%
\begin{Examples}
\begin{ExampleCode}
## Not run: 
result <- predict_oyster(survey, "ostrea_edulis")
tol    <- get_species_tolerances("ostrea_edulis")

# Project under all three UKCP18-aligned scenarios
proj <- project_suitability(result, tol)

# Mean suitability change by scenario
proj$summary

# High-risk cells: currently High but drops to Low under RCP8.5
vulnerable <- subset(proj$rcp85,
  result$suitability_class == "High" & suitability_class == "Low")

# Custom scenario (e.g. local downscaled projection)
proj2 <- project_suitability(result, tol, scenarios = NULL,
  custom_deltas = list(
    ukcp18_p95 = c(temperature = 3.2, salinity = -0.8)))

## End(Not run)
\end{ExampleCode}
\end{Examples}
\HeaderA{qc\_survey\_data}{Run automated quality control on raw survey data}{qc.Rul.survey.Rul.data}
%
\begin{Description}
Detects and flags suspect sensor readings before they enter the suitability
model. Three complementary checks are applied to each numeric column:
\begin{enumerate}

\item{} \strong{Range check} \bsl{}u2014 values outside the physically plausible range for that
variable (e.g. temperature > 35\bsl{}u00b0C in temperate coastal water, salinity42 PSU) are flagged as "range\_fail".
\item{} \strong{Statistical outlier} \bsl{}u2014 values beyond \code{iqr\_k} \bsl{}u00d7 IQR from the median
are flagged as \code{"outlier"} (default k = 3, Tukey's outer fence).
\item{} \strong{Temporal gradient} \bsl{}u2014 when a \code{datetime} column is present and data is
sorted chronologically, sequential differences exceeding the
\code{max\_gradient} threshold are flagged as \code{"gradient\_fail"} (instrument
spike or data entry error).

\end{enumerate}


Additionally, \strong{cross-variable sanity checks} flag physically implausible
combinations:
\begin{itemize}

\item{} Dissolved oxygen > 20 mg/L with temperature < 5\bsl{}u00b0C (likely sensor error)
\item{} Salinity < 5 PSU with temperature > 25\bsl{}u00b0C (likely freshwater intrusion
instrument error, not a real coastal reading)
\item{} Chlorophyll\_a > 50 \bsl{}u00b5g/L (extreme bloom or sensor fouling)

\end{itemize}


Flagged values are \strong{not removed} \bsl{}u2014 they are marked in a \AsIs{\texttt{qc\_flag\_<col>}}
column so the user can decide whether to replace with NA, correct, or
accept them. A summary \code{qc\_n\_flags} column counts total flags per row.
\end{Description}
%
\begin{Usage}
\begin{verbatim}
qc_survey_data(
  df,
  datetime_col = "datetime",
  iqr_k = 3,
  max_gradients = NULL,
  apply_flags = FALSE,
  verbose = TRUE
)
\end{verbatim}
\end{Usage}
%
\begin{Arguments}
\begin{ldescription}
\item[\code{df}] Dataframe of raw survey measurements.

\item[\code{datetime\_col}] Character or NULL. Name of the datetime column for
temporal gradient checks (default \code{"datetime"}). Set to NULL to skip.

\item[\code{iqr\_k}] Numeric. IQR multiplier for outlier detection (default 3.0;
use 1.5 for stricter QC of high-precision CTD data).

\item[\code{max\_gradients}] Named numeric vector. Maximum allowed absolute change
per unit time (seconds) per variable. Defaults are conservative for typical
towed or lowered deployments.

\item[\code{apply\_flags}] Logical. If TRUE, replace flagged values with NA in the
returned dataframe. If FALSE (default), flags are added but values
preserved for manual review.

\item[\code{verbose}] Logical. Print QC summary (default TRUE).
\end{ldescription}
\end{Arguments}
%
\begin{Value}
Input dataframe with additional columns:
\AsIs{\texttt{qc\_flag\_<varname>}} for each checked variable (\code{"ok"}, \code{"range\_fail"},
\code{"outlier"}, \code{"gradient\_fail"}, \code{"cross\_fail"}),
\code{qc\_n\_flags} (integer count of flags per row),
\code{qc\_status} (\code{"pass"}, \code{"review"}, \code{"fail"} \bsl{}u2014 fail = 3+ flags on one row).
\end{Value}
%
\begin{Examples}
\begin{ExampleCode}
## Not run: 
# Run QC before modelling
survey_qc <- qc_survey_data(survey, datetime_col = "timestamp")

# Inspect flagged rows
flagged <- subset(survey_qc, qc_status %in% c("review","fail"))
View(flagged)

# Apply flags (replace flagged values with NA) before scoring
survey_clean <- qc_survey_data(survey, apply_flags = TRUE)
result <- predict_oyster(survey_clean, "ostrea_edulis")

## End(Not run)
\end{ExampleCode}
\end{Examples}
\HeaderA{read\_generic\_csv}{Read a generic sensor CSV with flexible column mapping}{read.Rul.generic.Rul.csv}
%
\begin{Description}
Reads any sensor CSV (Lowrance BioBase, CTD/probe logger, bathymetric sonar,
sidescan export, etc.) and maps its columns to the standard oystermapR
variable names. Performs spatial averaging at a specified resolution.
\end{Description}
%
\begin{Usage}
\begin{verbatim}
read_generic_csv(
  file,
  col_map = NULL,
  spatial_res = 4L,
  min_obs = 3L,
  categorical_cols = c("sediment_type", "benthic_communities", "biotope",
    "fishing_intensity"),
  verbose = TRUE
)
\end{verbatim}
\end{Usage}
%
\begin{Arguments}
\begin{ldescription}
\item[\code{file}] Character. Path to the CSV file.

\item[\code{col\_map}] Named character vector mapping oystermapR variable names to
column names in the CSV. Only variables you want to extract need to be
listed. Example:

\begin{alltt}col_map = c(
  lat                 = "Latitude",
  lon                 = "Longitude",
  depth               = "Depth_m",
  substrate_hardness  = "Hardness",
  temperature         = "Temp_C"
)
\end{alltt}


If \code{col\_map} is \code{NULL} (default), the function attempts automatic column
name matching using built-in synonym lists (case-insensitive).

\item[\code{spatial\_res}] Integer. Decimal places for spatial binning (default \code{4}).

\item[\code{min\_obs}] Integer. Minimum observations per cell (default \code{3}).

\item[\code{categorical\_cols}] Character vector. Column names (after mapping) that
should be treated as categorical \bsl{}u2014 mode (most common value) is returned
instead of mean. Default: \code{c("sediment\_type", "benthic\_communities", "biotope", "fishing\_intensity")}.

\item[\code{verbose}] Logical. Print a summary after loading (default \code{TRUE}).
\end{ldescription}
\end{Arguments}
%
\begin{Value}
A spatially averaged dataframe using oystermapR standard column names.
\end{Value}
%
\begin{Examples}
\begin{ExampleCode}
## Not run: 
# Lowrance BioBase export (automatic column matching)
biobase <- read_generic_csv("biobase_export.csv")

# With explicit column mapping
probe <- read_generic_csv(
  "ctd_data.csv",
  col_map = c(lat="Lat", lon="Lon", temperature="Temp", salinity="Sal",
              dissolved_oxygen="DO_mgl")
)

## End(Not run)
\end{ExampleCode}
\end{Examples}
\HeaderA{read\_nortek\_adcp}{Read and process a Nortek Signature 500 ADCP CSV file}{read.Rul.nortek.Rul.adcp}
%
\begin{Description}
Reads a raw merged Nortek Signature 500 ADCP CSV, automatically detects
and excludes depth bins contaminated by acoustic sidelobe interference,
computes current velocity and estimated bed shear stress from the deepest
clean bin, and returns a spatially averaged dataframe ready to merge with
other sensor data via \code{\LinkA{merge\_sensor\_data()}{merge.Rul.sensor.Rul.data}}.

\strong{Sidelobe detection:} Each bin is tested independently. A bin is flagged
as contaminated if more than \code{sidelobe\_threshold} fraction of its speed
readings exceed \code{max\_plausible\_speed}. Once a bin is flagged, all deeper
bins are also flagged (contamination propagates downward).

\strong{Shear stress accuracy:} Bed shear stress (\code{tau = rho * Cd * U\_bed\textasciicircum{}2})
is most accurate when computed from near-bed velocity. If near-bed bins are
contaminated and only near-surface bins remain, a warning is issued and the
\code{shear\_stress\_quality} column is set to \code{"low"}. If two or more clean bins
exist, the deepest clean bin is used and quality is \code{"moderate"} or
\code{"good"}.
\end{Description}
%
\begin{Usage}
\begin{verbatim}
read_nortek_adcp(
  file,
  spatial_res = 4L,
  min_obs = 5L,
  max_plausible_speed = 1.5,
  sidelobe_threshold = 0.1,
  rho = 1025,
  Cd = 0.002,
  verbose = TRUE
)
\end{verbatim}
\end{Usage}
%
\begin{Arguments}
\begin{ldescription}
\item[\code{file}] Character. Path to the Nortek merged CSV file.

\item[\code{spatial\_res}] Integer. Decimal places for lat/lon binning (default \code{4},
\bsl{}u224811 m cells at 56\bsl{}u00b0N). Reduce to \code{3} (\textasciitilde{}111 m) for coarser surveys.

\item[\code{min\_obs}] Integer. Minimum number of raw ensembles per spatial cell
to include in output (default \code{5}). Cells with fewer observations are
dropped as unreliable.

\item[\code{max\_plausible\_speed}] Numeric. Speed in m/s above which a reading is
considered a sidelobe spike (default \code{1.5} m/s \bsl{}u2014 appropriate for sheltered
coastal/loch environments). Increase to \code{2.5} for open coast surveys.

\item[\code{sidelobe\_threshold}] Numeric. Fraction of readings in a bin that must
exceed \code{max\_plausible\_speed} before the bin is flagged as contaminated
(default \code{0.10}, i.e. 10\%).

\item[\code{rho}] Numeric. Seawater density in kg/m\bsl{}u00b3 (default \code{1025}).

\item[\code{Cd}] Numeric. Drag coefficient for bed shear stress calculation
(default \code{0.002}, typical for mixed sandy/rocky coastal seabed).

\item[\code{verbose}] Logical. Print processing summary (default \code{TRUE}).
\end{ldescription}
\end{Arguments}
%
\begin{Value}
A dataframe with one row per spatial cell containing:
\begin{itemize}

\item{} \code{lat}, \code{lon} \bsl{}u2014 cell centroid coordinates
\item{} \code{date} \bsl{}u2014 earliest observation date in cell (character, \code{"YYYY-MM-DD"})
\item{} \code{current\_velocity} \bsl{}u2014 mean speed from deepest clean bin (m/s)
\item{} \code{current\_velocity\_sd} \bsl{}u2014 standard deviation within cell (m/s)
\item{} \code{current\_velocity\_p95} \bsl{}u2014 95th percentile speed (m/s)
\item{} \code{shear\_stress} \bsl{}u2014 estimated bed shear stress (N/m\bsl{}u00b2)
\item{} \code{shear\_stress\_quality} \bsl{}u2014 \code{"good"}, \code{"moderate"}, or \code{"low"}
\item{} \code{n\_ensembles} \bsl{}u2014 raw ensembles averaged into this cell
\item{} \code{bins\_used} \bsl{}u2014 which velocity bins contributed (e.g. \code{"bin1"})
\item{} \code{bins\_excluded} \bsl{}u2014 which bins were removed as contaminated

\end{itemize}

\end{Value}
%
\begin{Examples}
\begin{ExampleCode}
## Not run: 
adcp <- read_nortek_adcp("S104456A008_AWE_Melfort_merged.csv")
print(adcp)

## End(Not run)
\end{ExampleCode}
\end{Examples}
\HeaderA{read\_sonar\_tif}{Read a bathymetric or sidescan GeoTIFF and convert to oystermapR variables}{read.Rul.sonar.Rul.tif}
%
\begin{Description}
Reads a single-band GeoTIFF raster using \code{terra} and returns a spatially
averaged dataframe suitable for \code{\LinkA{merge\_sensor\_data()}{merge.Rul.sensor.Rul.data}}.

Two modes are supported, selected via the \code{type} argument:

\strong{\code{"bathy"}} \bsl{}u2014 Bathymetric DEM (e.g. from Ping 3DSS or post-processed
soundings). Extracts:
\begin{itemize}

\item{} \code{depth} \bsl{}u2014 raster cell value (metres)
\item{} \code{slope} \bsl{}u2014 computed using \code{\LinkA{terra::terrain()}{terra::terrain()}}
\item{} \code{roughness} \bsl{}u2014 Terrain Ruggedness Index via \code{\LinkA{terra::terrain()}{terra::terrain()}}

\end{itemize}


\strong{\code{"sidescan"}} \bsl{}u2014 Sidescan backscatter mosaic (normalised 0\bsl{}u20131 or raw DN).
Extracts:
\begin{itemize}

\item{} \code{substrate\_hardness} \bsl{}u2014 backscatter intensity normalised to \LinkA{0, 1}{0, 1}.
High backscatter \bsl{}u2192 hard/coarse substrate (rock, gravel, shell);
low backscatter \bsl{}u2192 soft substrate (mud, fine sand).

\end{itemize}

\end{Description}
%
\begin{Usage}
\begin{verbatim}
read_sonar_tif(
  file,
  type = c("bathy", "sidescan"),
  spatial_res = 4L,
  band = 1L,
  nodata_val = NA,
  depth_positive = TRUE,
  verbose = TRUE
)
\end{verbatim}
\end{Usage}
%
\begin{Arguments}
\begin{ldescription}
\item[\code{file}] Character. Path to the GeoTIFF file (\code{.tif}).

\item[\code{type}] Character. Either \code{"bathy"} or \code{"sidescan"}.

\item[\code{spatial\_res}] Integer. Decimal places for output lat/lon grid
(default \code{4}). The raster is resampled to this resolution before
converting to a dataframe.

\item[\code{band}] Integer. Raster band to read (default \code{1}).

\item[\code{nodata\_val}] Numeric. Value to treat as no-data (default \code{NA}).
Override if your raster uses a sentinel (e.g. \code{-9999}).

\item[\code{depth\_positive}] Logical. For \code{type = "bathy"}: if \code{TRUE} (default),
raster values are already positive-downward. Set \code{FALSE} if stored as
negative elevation.

\item[\code{verbose}] Logical (default \code{TRUE}).
\end{ldescription}
\end{Arguments}
%
\begin{Value}
A dataframe with \code{lat}, \code{lon}, and the derived variable columns
ready for \code{\LinkA{merge\_sensor\_data()}{merge.Rul.sensor.Rul.data}}.
\end{Value}
%
\begin{Examples}
\begin{ExampleCode}
## Not run: 
bathy_df   <- read_sonar_tif("kames_bay_bathy.tif",    type = "bathy")
sidescan_df <- read_sonar_tif("kames_bay_sidescan.tif", type = "sidescan")
survey <- merge_sensor_data(bathy = bathy_df, sidescan = sidescan_df, adcp = adcp_df)

## End(Not run)
\end{ExampleCode}
\end{Examples}
\HeaderA{read\_soundings\_xyz}{Read a bathymetric XYZ point cloud and derive depth, slope and roughness}{read.Rul.soundings.Rul.xyz}
%
\begin{Description}
Reads a plain-text XYZ soundings file (comma or space delimited, optional
\AsIs{\texttt{\#}}-prefixed header), spatially averages onto a regular grid, and derives:
\begin{itemize}

\item{} \strong{\code{depth}} \bsl{}u2014 mean depth per cell (metres, positive downward)
\item{} \strong{\code{roughness}} \bsl{}u2014 rugosity index per cell, derived from intra-cell depth
variance using the surface-area approximation
\code{rugosity = sqrt(1 + (depth\_sd / cell\_size\_m)\textasciicircum{}2)}. Values near 1.0 are
flat; higher values indicate more complex relief.
\item{} \strong{\code{slope}} \bsl{}u2014 maximum downslope gradient in degrees, computed by finite
differences between neighbouring grid cells (Horn 1981 method).

\end{itemize}

\end{Description}
%
\begin{Usage}
\begin{verbatim}
read_soundings_xyz(
  file,
  lon_col = NULL,
  lat_col = NULL,
  depth_col = NULL,
  spatial_res = 4L,
  min_soundings = 10L,
  depth_positive = TRUE,
  verbose = TRUE
)
\end{verbatim}
\end{Usage}
%
\begin{Arguments}
\begin{ldescription}
\item[\code{file}] Character. Path to the \code{.xyz} (or \code{.csv}) soundings file.
Expected columns: longitude, latitude, depth (in that order, or named).
Lines beginning with \AsIs{\texttt{\#}} are treated as comments.

\item[\code{lon\_col}, \code{lat\_col}, \code{depth\_col}] Character or integer. Column names or
positions for longitude, latitude, and depth. If \code{NULL} (default), the
function tries common names automatically.

\item[\code{spatial\_res}] Integer. Decimal places for lat/lon binning (default \code{4},
\bsl{}u224811 m at 56\bsl{}u00b0N). The rugosity calculation uses the physical cell size
implied by this resolution.

\item[\code{min\_soundings}] Integer. Minimum soundings per cell to include in output
(default \code{10}). Cells with fewer soundings produce unreliable statistics.

\item[\code{depth\_positive}] Logical. If \code{TRUE} (default), depth values are already
positive-downward. Set to \code{FALSE} if your sonar records depth as negative
(elevation-style) \bsl{}u2014 values will be negated.

\item[\code{verbose}] Logical. Print processing summary (default \code{TRUE}).
\end{ldescription}
\end{Arguments}
%
\begin{Value}
A dataframe with columns \code{lat}, \code{lon}, \code{depth}, \code{slope},
\code{roughness}, and \code{n\_soundings} \bsl{}u2014 ready for \code{\LinkA{merge\_sensor\_data()}{merge.Rul.sensor.Rul.data}}.
\end{Value}
%
\begin{Examples}
\begin{ExampleCode}
## Not run: 
bathy <- read_soundings_xyz("kames_bay_soundings.xyz")
survey <- merge_sensor_data(adcp = adcp_data, bathy = bathy)

## End(Not run)
\end{ExampleCode}
\end{Examples}
\HeaderA{reset\_tolerance\_update}{Reset Bayesian tolerance updates for a species}{reset.Rul.tolerance.Rul.update}
%
\begin{Description}
Clears the session cache for a species, reverting to the built-in prior
parameters. Does not affect saved files (use file deletion for those).
\end{Description}
%
\begin{Usage}
\begin{verbatim}
reset_tolerance_update(species = "all", verbose = TRUE)
\end{verbatim}
\end{Usage}
%
\begin{Arguments}
\begin{ldescription}
\item[\code{species}] Character. Species key, or \code{"all"} to clear all cached updates.

\item[\code{verbose}] Logical. Default TRUE.
\end{ldescription}
\end{Arguments}
%
\begin{Value}
Invisibly NULL.
\end{Value}
%
\begin{Examples}
\begin{ExampleCode}
## Not run: 
reset_tolerance_update("ostrea_edulis")
reset_tolerance_update("all")

## End(Not run)
\end{ExampleCode}
\end{Examples}
\HeaderA{save\_tolerance\_update}{Save Bayesian tolerance updates to disk}{save.Rul.tolerance.Rul.update}
%
\begin{Description}
Serialises the current session cache for a species to an \code{.rds} file so
that updated parameters persist across R sessions. Load again with
\code{\LinkA{load\_tolerance\_update()}{load.Rul.tolerance.Rul.update}}.
\end{Description}
%
\begin{Usage}
\begin{verbatim}
save_tolerance_update(species = "all", path = "~/.oystermapR", verbose = TRUE)
\end{verbatim}
\end{Usage}
%
\begin{Arguments}
\begin{ldescription}
\item[\code{species}] Character. Species key or \code{"all"} to save all cached species.

\item[\code{path}] Character. Directory to save into (default: user cache directory
\AsIs{\texttt{\textasciitilde{}/.oystermapR/}}). Created if absent.

\item[\code{verbose}] Logical. Default TRUE.
\end{ldescription}
\end{Arguments}
%
\begin{Value}
Invisibly: path(s) to saved file(s).
\end{Value}
%
\begin{Examples}
\begin{ExampleCode}
## Not run: 
save_tolerance_update("ostrea_edulis")
save_tolerance_update("all", path = "data/bayes_updates/")

## End(Not run)
\end{ExampleCode}
\end{Examples}
\HeaderA{score\_anthropogenic\_disturbance}{Score anthropogenic disturbance at survey locations}{score.Rul.anthropogenic.Rul.disturbance}
%
\begin{Description}
Calculates a disturbance index \LinkA{0,1}{0,1} from bottom trawling intensity,
dredging activity, and anchor damage. The primary data input is the
ICES swept-area ratio (SAR), which can be supplied manually or fetched
live from the ICES VMS data portal.

This output is primarily relevant to the gear feasibility and economic
viability modules. Very high trawling intensity (SAR > 0.5/yr) constitutes
a hard gate for bottom culture and restoration reef deployment \bsl{}u2014 physical
gear deployment is not viable on actively trawled ground regardless of
ecological suitability.

\strong{Input options:}
\begin{enumerate}

\item{} \code{trawling\_data} \bsl{}u2014 ICES VMS SAR data, manually downloaded as CSV/dataframe.
\item{} \code{fetch\_live = TRUE} \bsl{}u2014 queries ICES VMS GeoServer endpoint.
\item{} Neither \bsl{}u2014 returns zero disturbance with a warning.

\end{enumerate}

\end{Description}
%
\begin{Usage}
\begin{verbatim}
score_anthropogenic_disturbance(
  result,
  trawling_data = NULL,
  anchor_data = NULL,
  dredge_areas = NULL,
  fetch_live = FALSE,
  match_radius_m = 5000,
  verbose = TRUE
)
\end{verbatim}
\end{Usage}
%
\begin{Arguments}
\begin{ldescription}
\item[\code{result}] Dataframe from \code{\LinkA{predict\_oyster()}{predict.Rul.oyster}}. Must contain \code{lat}, \code{lon}.

\item[\code{trawling\_data}] Dataframe of ICES VMS SAR data (see Details), or NULL.

\item[\code{anchor\_data}] Dataframe of anchoring density data, or NULL.

\item[\code{dredge\_areas}] Dataframe of dredging/extraction area centroids, or NULL.

\item[\code{fetch\_live}] Logical. Query ICES VMS GeoServer (default FALSE).

\item[\code{match\_radius\_m}] Numeric. Spatial matching radius in metres (default
5000 m, consistent with ICES c-square \textasciitilde{}0.05\bsl{}u00b0 resolution).

\item[\code{verbose}] Logical. Default TRUE.
\end{ldescription}
\end{Arguments}
%
\begin{Value}
\code{result} with additional columns:
\begin{itemize}

\item{} \code{disturbance\_sar}: swept-area ratio at or near each site (NA if no data)
\item{} \code{disturbance\_score} \LinkA{0,1}{0,1}: composite anthropogenic disturbance index
\item{} \code{disturbance\_class}: "Negligible" / "Low" / "Moderate" / "High" / "Active"
\item{} \code{disturbance\_gear\_gate}: logical \bsl{}u2014 TRUE if SAR exceeds gear deployment
threshold (used by \code{assess\_gear\_feasibility()})
\item{} \code{disturbance\_note}: management flag

\end{itemize}

\end{Value}
%
\begin{Section}{trawling\_data format}

A dataframe with columns:
\begin{itemize}

\item{} \code{lat}, \code{lon} \bsl{}u2014 centroid of c-square or fishing location
\item{} \code{sar} \bsl{}u2014 swept-area ratio (numeric; dimensionless ratio per year)
\item{} \code{gear\_type} (optional) \bsl{}u2014 "otter", "beam", "dredge", "seine" etc.
Dredge gear has higher impact per SAR unit.

\end{itemize}

\end{Section}
%
\begin{Section}{Additional disturbance inputs}

\begin{itemize}

\item{} \code{anchor\_data}: dataframe with \code{lat}, \code{lon}, \code{density} (vessel anchoring
density index) for recreational/commercial anchorage areas.
\item{} \code{dredge\_areas}: dataframe with \code{lat}, \code{lon} identifying licensed
aggregate extraction areas (all sites within the polygon receive score 1.0).

\end{itemize}

\end{Section}
%
\begin{Examples}
\begin{ExampleCode}
## Not run: 
# Manual ICES VMS data (downloaded from ICES data portal)
vms <- read.csv("ices_vms_sar_2022.csv")  # columns: lat, lon, sar, gear_type
result <- score_anthropogenic_disturbance(result, trawling_data = vms)

# Live ICES VMS fetch
result <- score_anthropogenic_disturbance(result, fetch_live = TRUE)

## End(Not run)
\end{ExampleCode}
\end{Examples}
\HeaderA{score\_disease\_risk}{Score locations for disease and parasite risk}{score.Rul.disease.Rul.risk}
%
\begin{Description}
Assesses the environmental risk of two key bivalve pathogens at each survey
location:
\begin{itemize}

\item{} \strong{Bonamia ostreae} (\emph{Ostrea edulis} only) \bsl{}u2014 an intracellular haplosporidian
parasite that is the primary constraint on flat oyster restoration in northern
Europe. Risk is temperature-driven: transmission peaks 15-20\bsl{}u00b0C and is
negligible below 8\bsl{}u00b0C. The function also accepts a \code{known\_sites} dataframe of
confirmed infected locations to apply a proximity multiplier.
\item{} \strong{OsHV-1 microvariant} (\emph{Magallana gigas}) \bsl{}u2014 a herpesvirus causing Pacific
Oyster Mortality Syndrome (POMS). Mortality events occur when seawater
exceeds 16\bsl{}u00b0C for sustained periods. High turbidity amplifies risk by
increasing larval stress.

\end{itemize}


Risk scores are returned on a 0-1 scale (0 = negligible, 1 = highest
environmental risk) and classified as Low / Moderate / High / Critical.
They are deliberately kept separate from the suitability score \bsl{}u2014 a High
suitability site with High disease risk is a meaningful regulatory red flag.
\end{Description}
%
\begin{Usage}
\begin{verbatim}
score_disease_risk(
  result,
  species = "ostrea_edulis",
  known_sites = NULL,
  temp_col = "temperature",
  salinity_col = "salinity",
  turbidity_col = "turbidity",
  verbose = TRUE
)
\end{verbatim}
\end{Usage}
%
\begin{Arguments}
\begin{ldescription}
\item[\code{result}] Dataframe from \code{\LinkA{predict\_oyster()}{predict.Rul.oyster}} with \code{lat}, \code{lon}, and
environmental variable columns.

\item[\code{species}] Character. \code{"ostrea\_edulis"} (Bonamia risk) or
\code{"magallana\_gigas"} (OsHV-1 risk).

\item[\code{known\_sites}] Dataframe or NULL. Known infected/positive sites with
\code{lat} and \code{lon} columns. When supplied, a proximity multiplier is applied.
For \emph{O. edulis} these are Bonamia-positive sites; for \emph{M. gigas} OsHV-1
outbreak sites.

\item[\code{temp\_col}] Character. Name of temperature column in \code{result}
(default \code{"temperature"}).

\item[\code{salinity\_col}] Character. Name of salinity column (default \code{"salinity"}).
Used for Bonamia only.

\item[\code{turbidity\_col}] Character. Name of turbidity column (default
\code{"turbidity"}). Used for OsHV-1 amplification only.

\item[\code{verbose}] Logical. Print risk summary (default TRUE).
\end{ldescription}
\end{Arguments}
%
\begin{Value}
Input dataframe with additional columns:
\code{disease\_risk\_score} \LinkA{0-1}{0.Rdash.1}, \code{disease\_risk\_class} (Low/Moderate/High/Critical),
\code{disease\_agent} (pathogen name), and \code{disease\_risk\_note} (plain-language
interpretation).
\end{Value}
%
\begin{Examples}
\begin{ExampleCode}
## Not run: 
result <- predict_oyster(survey, "ostrea_edulis")

# Basic risk scoring (temperature-driven only)
result <- score_disease_risk(result, "ostrea_edulis")

# With known Bonamia-positive site locations
infected <- data.frame(lat = c(55.8, 56.1), lon = c(-5.2, -5.4))
result <- score_disease_risk(result, "ostrea_edulis",
                              known_sites = infected)

# Sites with high ecological suitability but also high disease risk
flagged <- subset(result,
  suitability_class == "High" & disease_risk_class %in% c("High","Critical"))

## End(Not run)
\end{ExampleCode}
\end{Examples}
\HeaderA{score\_economic\_viability}{Score economic viability for aquaculture or restoration}{score.Rul.economic.Rul.viability}
%
\begin{Description}
Combines ecological suitability, gear feasibility, site accessibility, and
patch size into a composite \strong{economic viability index} \LinkA{0-1}{0.Rdash.1}. This is a
heuristic scoring \bsl{}u2014 not a financial model \bsl{}u2014 but it ranks sites consistently
and can be used to shortlist candidates for detailed business case
development.

The index weighs four components:
\begin{enumerate}

\item{} \strong{Ecological suitability} (40\%) \bsl{}u2014 from \code{\LinkA{predict\_oyster()}{predict.Rul.oyster}}
\item{} \strong{Gear feasibility} (25\%) \bsl{}u2014 best gear score from \code{\LinkA{assess\_gear\_feasibility()}{assess.Rul.gear.Rul.feasibility}}
\item{} \strong{Accessible patch area} (20\%) \bsl{}u2014 log-scaled area of connected suitable
habitat (from \code{\LinkA{analyse\_connectivity()}{analyse.Rul.connectivity}} if available)
\item{} \strong{Access score} (15\%) \bsl{}u2014 distance to nearest harbour/port (if
\code{harbours} table supplied) or flat 0.5 if unknown

\end{enumerate}

\end{Description}
%
\begin{Usage}
\begin{verbatim}
score_economic_viability(
  result,
  harbours = NULL,
  target = c("restoration", "aquaculture"),
  verbose = TRUE
)
\end{verbatim}
\end{Usage}
%
\begin{Arguments}
\begin{ldescription}
\item[\code{result}] Dataframe from \code{\LinkA{assess\_gear\_feasibility()}{assess.Rul.gear.Rul.feasibility}} (which itself
requires \code{\LinkA{predict\_oyster()}{predict.Rul.oyster}} output). Optionally also run
\code{\LinkA{analyse\_connectivity()}{analyse.Rul.connectivity}} first to enable area scoring.

\item[\code{harbours}] Dataframe or NULL. Known harbour/landing locations with
\code{lat} and \code{lon} columns. Used to score site accessibility.

\item[\code{target}] Character. \code{"restoration"} (default) weights ecological
suitability and connectivity higher. \code{"aquaculture"} weights gear
feasibility and economic access higher.

\item[\code{verbose}] Logical. Print viability summary (default TRUE).
\end{ldescription}
\end{Arguments}
%
\begin{Value}
Input dataframe with \code{viability\_score} \LinkA{0-1}{0.Rdash.1},
\code{viability\_class} (Poor/Fair/Good/Excellent), and \code{viability\_notes}.
\end{Value}
%
\begin{Examples}
\begin{ExampleCode}
## Not run: 
result <- predict_oyster(survey, "ostrea_edulis")
result <- assess_gear_feasibility(result)
result <- analyse_connectivity(result)

harbours <- data.frame(
  name = c("Tarbert","Portavadie"),
  lat  = c(55.865, 55.875),
  lon  = c(-5.425, -5.300)
)

result <- score_economic_viability(result, harbours = harbours,
                                    target = "aquaculture")

# Best aquaculture candidates
subset(result, viability_class %in% c("Good","Excellent"))

## End(Not run)
\end{ExampleCode}
\end{Examples}
\HeaderA{score\_hab\_risk}{Score harmful algal bloom risk at survey locations}{score.Rul.hab.Rul.risk}
%
\begin{Description}
Produces a HAB risk index (0 = negligible, 1 = severe) based on historical
bloom event frequency and, where available, nutrient and stratification data.

\strong{Input options (in priority order):}
\begin{enumerate}

\item{} \code{hab\_data} \bsl{}u2014 manually supplied dataframe of historical HAB events
(see Details).
\item{} \code{fetch\_live = TRUE} \bsl{}u2014 queries the ICES HAB database for event records
within the survey extent and a specified date range. Requires internet
access and the \code{httr} package.
\item{} Neither \bsl{}u2014 returns a fixed background-level risk (0.1) with a warning.

\end{enumerate}

\end{Description}
%
\begin{Usage}
\begin{verbatim}
score_hab_risk(
  result,
  hab_data = NULL,
  fetch_live = FALSE,
  date_range = NULL,
  match_radius_m = 5000,
  species = "ostrea_edulis",
  verbose = TRUE
)
\end{verbatim}
\end{Usage}
%
\begin{Arguments}
\begin{ldescription}
\item[\code{result}] Dataframe from \code{\LinkA{predict\_oyster()}{predict.Rul.oyster}}. Must contain \code{lat}, \code{lon}.
Optional: \code{din\_ug\_l} (dissolved inorganic nitrogen \bsl{}u00b5g/L) and
\code{temp\_surface\_c} / \code{temp\_bottom\_c} for stratification index.

\item[\code{hab\_data}] Dataframe of historical HAB events (see Details), or NULL.

\item[\code{fetch\_live}] Logical. Query ICES HAB database (default FALSE).

\item[\code{date\_range}] Character vector length 2 \code{c("YYYY-MM-DD","YYYY-MM-DD")}.
Date range for live fetch or for filtering manually supplied \code{hab\_data}.
Default: last 10 years.

\item[\code{match\_radius\_m}] Numeric. Radius in metres for event aggregation
(default 5000 m = 5 km, consistent with ICES monitoring station spacing).

\item[\code{species}] Character. Target species \bsl{}u2014 affects toxin severity weights
(\code{"ostrea\_edulis"} or \code{"magallana\_gigas"}).

\item[\code{verbose}] Logical. Default TRUE.
\end{ldescription}
\end{Arguments}
%
\begin{Value}
\code{result} with additional columns:
\begin{itemize}

\item{} \code{hab\_risk} \LinkA{0,1}{0,1}: composite risk index
\item{} \code{hab\_risk\_class}: "Low" / "Moderate" / "High" / "Critical"
\item{} \code{hab\_closure\_days\_per\_year}: estimated days/year with shellfish closures
\item{} \code{hab\_dominant\_toxin}: most frequently recorded toxin class at the site
\item{} \code{hab\_risk\_note}: management flag

\end{itemize}

\end{Value}
%
\begin{Section}{hab\_data format}

A dataframe with columns:
\begin{itemize}

\item{} \code{lat}, \code{lon} \bsl{}u2014 event coordinates (decimal degrees)
\item{} \code{date} \bsl{}u2014 event date (Date or character "YYYY-MM-DD")
\item{} \code{genus} (optional) \bsl{}u2014 bloom genus (e.g. "Alexandrium", "Pseudo-nitzschia")
\item{} \code{toxin} (optional) \bsl{}u2014 toxin class (PSP, ASP, DSP, AZP) \bsl{}u2014 affects severity weight
\item{} \code{closure\_days} (optional) \bsl{}u2014 number of days the area was closed; default 1
per event if absent

\end{itemize}

\end{Section}
%
\begin{Examples}
\begin{ExampleCode}
## Not run: 
# Manual data
hab <- data.frame(lat=52.1, lon=-4.5, date="2022-07-15",
                  genus="Alexandrium", toxin="PSP", closure_days=14)
result <- score_hab_risk(result, hab_data=hab, species="ostrea_edulis")

# Live ICES fetch
result <- score_hab_risk(result, fetch_live=TRUE,
                         date_range=c("2015-01-01","2024-12-31"))

## End(Not run)
\end{ExampleCode}
\end{Examples}
\HeaderA{score\_larval\_connectivity}{Score larval dispersal connectivity at survey locations}{score.Rul.larval.Rul.connectivity}
%
\begin{Description}
Estimates the larval connectivity of each survey location using one or both
of two approaches:

\strong{Route 1 \bsl{}u2014 Built-in gap-threshold union-find (no external data needed):}
Patches above \code{min\_suitability} that lie within the species' effective
dispersal kernel distance are grouped into dispersal clusters using
union-find. Within each cluster, a connectivity score is derived from the
number, quality, and distance-weighted contribution of potential larval
sources (Gaussian decay kernel, \bsl{}u03c3 = dispersal\_km / 2).

\strong{Route 2 \bsl{}u2014 External connectivity matrix (particle tracking / literature):}
When \code{connectivity\_matrix} is supplied (a dataframe of source \bsl{}u2192 destination
pairs with weights), those weights replace the Gaussian kernel for matched
pairs. Rows not covered by the matrix fall back to Route 1. The matrix can
come from:
\begin{itemize}

\item{} \strong{OpenDrift} particle tracking simulations (export as source/destination
settlement density CSV)
\item{} \strong{FVCOM / ROMS} model connectivity matrices
\item{} \strong{Published empirical matrices} (e.g. Robins et al. 2017)

\end{itemize}

\end{Description}
%
\begin{Usage}
\begin{verbatim}
score_larval_connectivity(
  result,
  species = NULL,
  dispersal_km = NULL,
  pld_days = NULL,
  tidal_excursion_km = 5,
  min_suitability = 0.4,
  connectivity_matrix = NULL,
  matrix_match_radius_deg = 0.1,
  verbose = TRUE
)
\end{verbatim}
\end{Usage}
%
\begin{Arguments}
\begin{ldescription}
\item[\code{result}] Dataframe from \code{\LinkA{predict\_oyster()}{predict.Rul.oyster}} with \code{lat}, \code{lon},
\code{suitability} columns.

\item[\code{species}] Character. Species key (e.g. \code{"ostrea\_edulis"}). Used to
look up pelagic larval duration for dispersal kernel. Required if
\code{dispersal\_km} is not supplied.

\item[\code{dispersal\_km}] Numeric. Override for the effective dispersal kernel
radius in km. When supplied, \code{species} PLD look-up is skipped.

\item[\code{pld\_days}] Numeric. Override for the mean pelagic larval duration in
days. Takes precedence over species lookup; ignored if \code{dispersal\_km} is
supplied.

\item[\code{tidal\_excursion\_km}] Numeric. Mean daily tidal dispersal distance in km
(default 5 km/day \bsl{}u2248 mean current 0.06 m/s \bsl{}u00d7 24 h). Increase for highly
tidal or estuarine sites (8\bsl{}u201315 km/day). Ignored if \code{dispersal\_km} is
supplied.

\item[\code{min\_suitability}] Numeric. Minimum suitability score for a patch to
qualify as a potential larval source or sink (default 0.40). Unsuitable
patches still receive a connectivity score reflecting their proximity to
sources.

\item[\code{connectivity\_matrix}] Dataframe or NULL. External connectivity weights
(Route 2). Must contain columns:
\begin{itemize}

\item{} \code{source\_lat}, \code{source\_lon}: coordinates of the larval source
\item{} \code{dest\_lat}, \code{dest\_lon}: coordinates of the destination
\item{} \code{weight}: connectivity weight \LinkA{0, 1}{0, 1} (e.g. settlement probability or
proportional larval flux). Higher = stronger connection.
OpenDrift export tip: \AsIs{\texttt{\bsl{}dontrun\{write.csv(connectivity\_matrix, "cm.csv")\}}}.

\end{itemize}


\item[\code{matrix\_match\_radius\_deg}] Numeric. Spatial matching tolerance in decimal
degrees for linking matrix entries to result rows (default 0.10 \bsl{}u2248 10 km).

\item[\code{verbose}] Logical. Print dispersal parameters and connectivity summary
(default TRUE).
\end{ldescription}
\end{Arguments}
%
\begin{Value}
\code{result} with additional columns:
\begin{itemize}

\item{} \code{larval\_dispersal\_km}: effective kernel distance used
\item{} \code{larval\_pld\_days}: PLD used (species default or override)
\item{} \code{larval\_type}: "lecithotrophic" or "planktotrophic"
\item{} \code{n\_larval\_sources}: suitable patches within dispersal range
\item{} \code{source\_quality\_score} \LinkA{0,1}{0,1}: Gaussian-weighted quality of nearby sources
\item{} \code{larval\_cluster\_id}: dispersal network cluster ID (integer; NA if
below min\_suitability threshold and no sources nearby)
\item{} \code{larval\_cluster\_size}: number of suitable patches in dispersal network
\item{} \code{nearest\_source\_km}: km to nearest suitable source patch
\item{} \code{larval\_connectivity\_score} \LinkA{0,1}{0,1}: composite dispersal connectivity score
\item{} \code{larval\_connectivity\_class}: "Highly connected" / "Connected" /
"Low connectivity" / "Isolated"
\item{} \code{larval\_connectivity\_note}: plain-language ecological interpretation
\item{} \code{larval\_source}: "union\_find", "matrix", or "matrix+union\_find"

\end{itemize}

\end{Value}
%
\begin{Section}{Dispersal distance}

If \code{dispersal\_km} is not supplied, the effective kernel distance is derived
from the species' mean pelagic larval duration (PLD) multiplied by
\code{tidal\_excursion\_km} (mean daily tidal dispersal distance):

\begin{alltt}dispersal_km = PLD_mean_days \bsl{}u00d7 tidal_excursion_km
\end{alltt}


PLD values are literature-derived:
\Tabular{llll}{
Species & PLD & Larval type & Default dispersal \\{}
\emph{O. edulis} & 1\bsl{}u20137 days & lecithotrophic & \textasciitilde{}20 km \\{}
\emph{M. gigas} & 14\bsl{}u201328 days & planktotrophic & \textasciitilde{}105 km \\{}
\emph{C. angulata} & 14\bsl{}u201321 days & planktotrophic & \textasciitilde{}85 km \\{}
\emph{O. stentina} & 2\bsl{}u20137 days & lecithotrophic & \textasciitilde{}20 km \\{}
\emph{O. lurida} & 3\bsl{}u201310 days & lecithotrophic & \textasciitilde{}30 km \\{}
}
\end{Section}
%
\begin{References}
Cowen R.K. \& Sponaugle S. (2009) Annual Review of Marine Science 1:443\bsl{}u2013466.
Robins P.E. et al. (2017) J Applied Ecology 54:1699\bsl{}u20131710. \Rhref{https://doi.org/10.1111/1365-2664.12854}{doi:10.1111\slash{}1365\-2664.12854}
Siegel D.A. et al. (2008) PNAS 105:8974\bsl{}u20138979.
Woolmer A.P. et al. (2009) J Shellfish Research 28:107\bsl{}u2013116.
Trimble A.C. et al. (2009) J Shellfish Research 28:43\bsl{}u201353.
\end{References}
%
\begin{Examples}
\begin{ExampleCode}
## Not run: 
result <- predict_oyster(survey, "ostrea_edulis")

# Route 1 only \u2014 built-in dispersal kernel
result <- score_larval_connectivity(result, species = "ostrea_edulis")

# Route 1 with custom dispersal distance (e.g. strong tidal excursion)
result <- score_larval_connectivity(result, species = "ostrea_edulis",
                                    tidal_excursion_km = 10)

# Route 1 with manual dispersal override
result <- score_larval_connectivity(result, dispersal_km = 8)

# Route 2 \u2014 external connectivity matrix from OpenDrift
cm <- read.csv("opendrift_connectivity.csv")
# Required columns: source_lat, source_lon, dest_lat, dest_lon, weight
result <- score_larval_connectivity(result,
           species             = "ostrea_edulis",
           connectivity_matrix = cm)

# Inspect isolated patches \u2014 need artificial seeding
isolated <- subset(result,
  larval_connectivity_class == "Isolated" & suitability_class == "High")

# Strong candidates: high suitability AND high connectivity
targets <- subset(result,
  suitability_class == "High" & larval_connectivity_class == "Highly connected")

## End(Not run)
\end{ExampleCode}
\end{Examples}
\HeaderA{score\_locations}{Score all locations in a dataframe}{score.Rul.locations}
%
\begin{Description}
Applies \code{\LinkA{.score\_row()}{.score.Rul.row}} across every non-excluded row. Excluded rows
receive \code{suitability = 0}. Adds per-variable score columns and a
\code{limiting\_factors} column identifying the variables most constraining
the score at each location.
\end{Description}
%
\begin{Usage}
\begin{verbatim}
score_locations(df, tolerances, verbose = FALSE)
\end{verbatim}
\end{Usage}
%
\begin{Arguments}
\begin{ldescription}
\item[\code{df}] A dataframe processed by \code{\LinkA{check\_exclusions()}{check.Rul.exclusions}}.

\item[\code{tolerances}] Species tolerance list from \code{\LinkA{get\_species\_tolerances()}{get.Rul.species.Rul.tolerances}}.

\item[\code{verbose}] Logical. Print per-variable scoring summary (default \code{FALSE}).
\end{ldescription}
\end{Arguments}
%
\begin{Value}
The input dataframe with additional columns:
\begin{itemize}

\item{} \code{suitability}: weighted score \LinkA{0, 1}{0, 1}
\item{} \code{suitability\_class}: "High" / "Moderate" / "Low" / "Very Low" / "Excluded"
\item{} \code{limiting\_factors}: comma-separated names of the variables most limiting
the score (NA if all variables score \bsl{}u2265 0.65)
\item{} \AsIs{\texttt{score\_<variable>}}: one column per scored variable present in the data

\end{itemize}

\end{Value}
\HeaderA{score\_predation\_risk}{Score predation and bioturbation pressure at survey locations}{score.Rul.predation.Rul.risk}
%
\begin{Description}
Produces a predation risk index (0 = negligible, 1 = severe) based on
predator occurrence data, substrate rugosity, and depth. The output is
intended as an \strong{overlay} on ecological suitability \bsl{}u2014 high predation risk
does not reduce the habitat suitability score directly but is flagged in
a separate column for management planning (e.g. cage exclusion, deep
subtidal placement to avoid intertidal drills).

\strong{Input options (in priority order):}
\begin{enumerate}

\item{} \code{predator\_data} \bsl{}u2014 manually supplied dataframe with predator records
(see Details).
\item{} \code{fetch\_live = TRUE} \bsl{}u2014 queries EMODnet Biology for \emph{Asterias rubens} and
\emph{Carcinus maenas} occurrence records within the survey extent. Requires
internet access and the \code{httr} package.
\item{} Neither \bsl{}u2014 returns a depth-proxy-only risk score with a warning.

\end{enumerate}

\end{Description}
%
\begin{Usage}
\begin{verbatim}
score_predation_risk(
  result,
  predator_data = NULL,
  fetch_live = FALSE,
  match_radius_m = 1000,
  species = "ostrea_edulis",
  verbose = TRUE
)
\end{verbatim}
\end{Usage}
%
\begin{Arguments}
\begin{ldescription}
\item[\code{result}] Dataframe from \code{\LinkA{predict\_oyster()}{predict.Rul.oyster}} with \code{lat}, \code{lon}, and
optionally \code{depth\_m} and \code{substrate} columns.

\item[\code{predator\_data}] Dataframe of predator records (see Details), or NULL.

\item[\code{fetch\_live}] Logical. Query EMODnet Biology API (default FALSE).

\item[\code{match\_radius\_m}] Numeric. Radius in metres within which predator
records are aggregated per survey cell (default 1000 m).

\item[\code{species}] Character. Target oyster species \bsl{}u2014 affects which predators
are most relevant (\code{"ostrea\_edulis"} or \code{"magallana\_gigas"}).

\item[\code{verbose}] Logical. Default TRUE.
\end{ldescription}
\end{Arguments}
%
\begin{Value}
\code{result} with additional columns:
\begin{itemize}

\item{} \code{predation\_risk} \LinkA{0,1}{0,1}: composite risk index
\item{} \code{predation\_risk\_class}: "Low" / "Moderate" / "High" / "Severe"
\item{} \code{predation\_risk\_note}: text flag for management planning
\item{} \code{n\_predator\_records}: number of predator records within match\_radius\_m

\end{itemize}

\end{Value}
%
\begin{Section}{predator\_data format}

A dataframe with columns:
\begin{itemize}

\item{} \code{lat}, \code{lon} \bsl{}u2014 coordinates
\item{} \code{species} \bsl{}u2014 species name or AphiaID (character)
\item{} \code{density} (optional) \bsl{}u2014 relative abundance / density index; if absent,
each record is treated as presence-only (density = 1)
\item{} \code{date} (optional) \bsl{}u2014 survey date (used to filter to relevant season)

\end{itemize}

\end{Section}
%
\begin{Examples}
\begin{ExampleCode}
## Not run: 
# Manual data
pred <- data.frame(lat=51.5, lon=-4.1, species="Asterias rubens", density=3)
result <- score_predation_risk(result, predator_data=pred, species="ostrea_edulis")

# Live EMODnet fetch
result <- score_predation_risk(result, fetch_live=TRUE)

# Depth-proxy only (no predator data)
result <- score_predation_risk(result)

## End(Not run)
\end{ExampleCode}
\end{Examples}
\HeaderA{score\_sediment\_stability}{Score sediment stability and mobility at survey locations}{score.Rul.sediment.Rul.stability}
%
\begin{Description}
Calculates the Shields parameter and a sediment stability score \LinkA{0,1}{0,1} for
each survey location. High sediment mobility (\bsl{}u03b8 >> \bsl{}u03b8\_cr) indicates that
the seabed is frequently in motion, which is unfavourable for oyster
settlement, juvenile survival, and restoration reef persistence.
\end{Description}
%
\begin{Usage}
\begin{verbatim}
score_sediment_stability(
  result,
  current_col = "current_ms",
  wave_hs_col = "wave_hs_m",
  depth_col = "depth_m",
  substrate_col = "substrate",
  d50_col = "d50_mm",
  wave_period_s = 8,
  drag_coef = 0.003,
  verbose = TRUE
)
\end{verbatim}
\end{Usage}
%
\begin{Arguments}
\begin{ldescription}
\item[\code{result}] Dataframe from \code{\LinkA{predict\_oyster()}{predict.Rul.oyster}}. Must contain \code{lat}, \code{lon}.
Key optional columns (column names configurable via arguments):
\begin{itemize}

\item{} \code{current\_ms}: depth-averaged current speed (m/s)
\item{} \code{wave\_hs\_m}: significant wave height (m) \bsl{}u2014 from \code{score\_wave\_exposure()}
or direct measurement
\item{} \code{depth\_m}: water depth (m) \bsl{}u2014 needed for wave orbital velocity
\item{} \code{substrate}: substrate classification string \bsl{}u2014 used to estimate d50
\item{} \code{d50\_mm}: median grain diameter (mm) \bsl{}u2014 overrides substrate-derived d50

\end{itemize}


\item[\code{current\_col}] Character. Column name for depth-averaged current speed
(m/s). Default \code{"current\_ms"}.

\item[\code{wave\_hs\_col}] Character. Column name for significant wave height (m).
Default \code{"wave\_hs\_m"}.

\item[\code{depth\_col}] Character. Column name for depth (m). Default \code{"depth\_m"}.

\item[\code{substrate\_col}] Character. Column name for substrate type. Default
\code{"substrate"}.

\item[\code{d50\_col}] Character. Column name for measured median grain size (mm).
Default \code{"d50\_mm"}. If present, overrides substrate-derived estimate.

\item[\code{wave\_period\_s}] Numeric. Peak wave period in seconds for orbital
velocity calculation. Default 8 s (typical NW European sea state).

\item[\code{drag\_coef}] Numeric. Quadratic bed drag coefficient C\_f. Default
0.003 (smooth mixed seabed; increase to 0.005-0.010 for rough rock/reef).

\item[\code{verbose}] Logical. Default TRUE.
\end{ldescription}
\end{Arguments}
%
\begin{Value}
\code{result} with additional columns:
\begin{itemize}

\item{} \code{shields\_parameter}: Shields number \bsl{}u03b8 at each site
\item{} \code{shields\_critical}: critical Shields number \bsl{}u03b8\_cr for the substrate d50
\item{} \code{mobility\_ratio}: \bsl{}u03b8 / \bsl{}u03b8\_cr (> 1.0 = mobile; < 1.0 = stable)
\item{} \code{d50\_mm\_estimated}: grain size used (mm) \bsl{}u2014 measured or substrate-derived
\item{} \code{sediment\_stability\_score} \LinkA{0,1}{0,1}: 1 = fully stable, 0 = highly mobile
\item{} \code{sediment\_mobility\_class}: "Stable" / "Marginally stable" / "Mobile" / "Highly mobile"
\item{} \code{sediment\_stability\_note}: management flag

\end{itemize}

\end{Value}
%
\begin{Section}{Inputs required}

At minimum, the function needs current velocity to estimate bed shear
stress. Grain size improves accuracy but can be estimated from the
substrate type column if not directly measured.
\end{Section}
%
\begin{Examples}
\begin{ExampleCode}
## Not run: 
# With current velocity and substrate type columns
result <- score_sediment_stability(result, current_col="current_ms", substrate_col="substrate")

# After score_wave_exposure() \u2014 wave_hs_m column already present
result <- score_wave_exposure(result, fetch_km=15)
result <- score_sediment_stability(result)

# With measured grain size
result$d50_mm <- c(0.25, 2.5, 15.0, 0.08)
result <- score_sediment_stability(result)

## End(Not run)
\end{ExampleCode}
\end{Examples}
\HeaderA{score\_settlement}{Score locations for spat settlement suitability}{score.Rul.settlement}
%
\begin{Description}
Assesses whether survey locations meet the conditions required for bivalve
larval settlement and metamorphosis. Settlement scoring uses a separate,
tighter tolerance specification than adult habitat scoring \bsl{}u2014 larvae are more
sensitive to turbidity, require minimum current flow for delivery, and need
hard substrate for attachment.

The output can be compared directly against the adult suitability score from
\code{\LinkA{predict\_oyster()}{predict.Rul.oyster}} to distinguish "good adult habitat" from "likely natural
recruitment site."
\end{Description}
%
\begin{Usage}
\begin{verbatim}
score_settlement(survey, species = "ostrea_edulis", verbose = TRUE)
\end{verbatim}
\end{Usage}
%
\begin{Arguments}
\begin{ldescription}
\item[\code{survey}] Dataframe of survey measurements. Same format as \code{\LinkA{predict\_oyster()}{predict.Rul.oyster}}.

\item[\code{species}] Character. One of \code{"ostrea\_edulis"} or \code{"magallana\_gigas"}.

\item[\code{verbose}] Logical. Print summary (default TRUE).
\end{ldescription}
\end{Arguments}
%
\begin{Value}
Dataframe with \code{settlement\_suitability}, \code{settlement\_class}, and
per-variable \AsIs{\texttt{settle\_score\_*}} columns. Combine with adult result via
\code{cbind()} or merge on \code{lat}/\code{lon}.
\end{Value}
%
\begin{Examples}
\begin{ExampleCode}
## Not run: 
adult_result     <- predict_oyster(survey, "ostrea_edulis")
settlement_result <- score_settlement(survey, "ostrea_edulis")

# Identify cells suitable for both adult survival AND natural recruitment
combined <- merge(adult_result, settlement_result[, c("lat","lon",
              "settlement_suitability","settlement_class")],
              by = c("lat","lon"))
combined$dual_suitable <- combined$suitability >= 0.6 &
                          combined$settlement_suitability >= 0.6

## End(Not run)
\end{ExampleCode}
\end{Examples}
\HeaderA{score\_wave\_exposure}{Score wave exposure from fetch or measured wave height}{score.Rul.wave.Rul.exposure}
%
\begin{Description}
Produces a wave exposure score \LinkA{0,1}{0,1} and significant wave height estimate
for each survey location. High wave exposure reduces gear feasibility and
can destabilise restoration reef structures and spat bags.

\strong{Input options:}
\begin{itemize}

\item{} \code{fetch\_km}: effective fetch in kilometres (distance to nearest land or
obstruction in the prevailing wind direction). Most useful for estuarine,
sea loch, or enclosed bay sites. Can be measured from GIS (e.g. in QGIS
using the Fetch tool) or estimated from chart inspection.
\item{} \code{wave\_height\_m}: directly measured or modelled significant wave height
(Hs). If supplied, \code{fetch\_km} is ignored for Hs computation.
\item{} Both absent: wave exposure is estimated from depth alone (coarse proxy;
a warning is issued).

\end{itemize}

\end{Description}
%
\begin{Usage}
\begin{verbatim}
score_wave_exposure(
  result,
  fetch_col = NULL,
  wave_height_col = NULL,
  fetch_km = NULL,
  wave_height_m = NULL,
  wind_speed_ms = 12,
  depth_limit = TRUE,
  verbose = TRUE
)
\end{verbatim}
\end{Usage}
%
\begin{Arguments}
\begin{ldescription}
\item[\code{result}] Dataframe from \code{\LinkA{predict\_oyster()}{predict.Rul.oyster}}. Must contain \code{lat}, \code{lon}.
Optional columns used if present: \code{depth\_m}, \code{fetch\_km}, \code{wave\_height\_m}.

\item[\code{fetch\_col}] Character. Name of a column in \code{result} containing fetch
values in km. Alternative to supplying a scalar \code{fetch\_km}.

\item[\code{wave\_height\_col}] Character. Name of a column in \code{result} containing
measured/modelled Hs in metres.

\item[\code{fetch\_km}] Numeric scalar. Uniform fetch in km applied to all rows
(overridden by \code{fetch\_col} if both supplied).

\item[\code{wave\_height\_m}] Numeric scalar. Uniform Hs in metres applied to all
rows (overridden by \code{wave\_height\_col} if both supplied).

\item[\code{wind\_speed\_ms}] Numeric. Design wind speed in m/s used for JONSWAP
calculation (default 12 m/s = moderate gale, Beaufort 6, typical design
condition for coastal aquaculture siting).

\item[\code{depth\_limit}] Logical. Apply depth-limited wave breaking cap
(Hs\_max = 0.6 * depth\_m). Default TRUE.

\item[\code{verbose}] Logical. Default TRUE.
\end{ldescription}
\end{Arguments}
%
\begin{Value}
\code{result} with additional columns:
\begin{itemize}

\item{} \code{wave\_hs\_m}: significant wave height (m) \bsl{}u2014 computed or supplied
\item{} \code{wave\_exposure\_score} \LinkA{0,1}{0,1}: 0 = fully sheltered, 1 = severely exposed
\item{} \code{wave\_exposure\_class}: "Sheltered" / "Moderate" / "Exposed" / "Severe"
\item{} \code{wave\_exposure\_note}: gear implication flag
\item{} \code{wave\_source}: method used ("measured", "jonswap\_fetch", "depth\_proxy")

\end{itemize}

\end{Value}
%
\begin{Examples}
\begin{ExampleCode}
## Not run: 
# Fetch column in result
result$fetch_km <- c(2.5, 8.0, 25.0, 45.0)
result <- score_wave_exposure(result, fetch_col = "fetch_km")

# Uniform fetch for a sheltered sea loch
result <- score_wave_exposure(result, fetch_km = 3.5)

# Measured wave height column
result <- score_wave_exposure(result, wave_height_col = "hs_m")

# Depth proxy only (coarse)
result <- score_wave_exposure(result)

## End(Not run)
\end{ExampleCode}
\end{Examples}
\HeaderA{sensitivity\_analysis}{Partial dependence / sensitivity analysis for a scored variable}{sensitivity.Rul.analysis}
%
\begin{Description}
Computes the marginal response curve for a single environmental variable:
suitability is predicted at a grid of values for the focal variable, while
all other variables are held at their observed median (or a specified
background value). This reveals the shape of the scoring function as
actually applied to a real dataset.

Partial dependence is a standard diagnostic for SDM publication (Elith et al.
2008; Zurell et al. 2020 ODMAP protocol). Use it to verify that predicted
responses match ecological expectations before submission.
\end{Description}
%
\begin{Usage}
\begin{verbatim}
sensitivity_analysis(
  predicted,
  species,
  variable,
  n_steps = 100L,
  background = NULL,
  season = NULL,
  verbose = TRUE
)
\end{verbatim}
\end{Usage}
%
\begin{Arguments}
\begin{ldescription}
\item[\code{predicted}] Dataframe from \code{\LinkA{predict\_oyster()}{predict.Rul.oyster}} containing \code{lat}, \code{lon},
\code{suitability}, and per-variable score/weight columns.

\item[\code{species}] Character. Species key (e.g. \code{"ostrea\_edulis"}). Used to
retrieve tolerance parameters for the scoring function.

\item[\code{variable}] Character. Name of the focal variable to vary (e.g.
\code{"temperature"}, \code{"salinity"}, \code{"depth"}).

\item[\code{n\_steps}] Integer. Number of evenly spaced values across the variable's
biological range (default 100).

\item[\code{background}] Named list. Fixed values for all other variables. If NULL
(default), uses column medians from \code{predicted}.

\item[\code{season}] Character or NULL. Season to apply for seasonal variables
(e.g. \code{"summer"}). NULL = no seasonal override.

\item[\code{verbose}] Logical. Default TRUE.
\end{ldescription}
\end{Arguments}
%
\begin{Value}
A dataframe with columns:
\begin{itemize}

\item{} \code{x}: focal variable value
\item{} \code{suitability}: predicted suitability at that value
\item{} \code{variable}: focal variable name
\item{} \code{species}: species key

\end{itemize}


Suitable for plotting with \code{ggplot2} or base R \code{plot()}.
\end{Value}
%
\begin{References}
Elith et al. (2008) J Animal Ecology 77:802-813.
Zurell et al. (2020) Ecography 43:1261-1277.
\end{References}
%
\begin{Examples}
\begin{ExampleCode}
## Not run: 
result <- predict_oyster(survey, "ostrea_edulis")

# Temperature response curve
temp_pd <- sensitivity_analysis(result, "ostrea_edulis", "temperature")
plot(temp_pd$x, temp_pd$suitability, type="l",
     xlab="Temperature (\u00b0C)", ylab="Suitability",
     main="Partial dependence: temperature")

# Salinity response (summer)
sal_pd <- sensitivity_analysis(result, "ostrea_edulis", "salinity",
                               season = "summer")

# ggplot2 version
library(ggplot2)
ggplot(temp_pd, aes(x, suitability)) +
  geom_line(colour="steelblue", linewidth=1.2) +
  geom_ribbon(aes(ymin=0, ymax=suitability), alpha=0.2, fill="steelblue") +
  labs(x="Temperature (\u00b0C)", y="Suitability [0-1]") +
  theme_minimal()

## End(Not run)
\end{ExampleCode}
\end{Examples}
\HeaderA{smooth\_suitability}{Apply Gaussian kernel smoothing to suitability scores}{smooth.Rul.suitability}
%
\begin{Description}
Individual CTD or ADCP casts are point measurements subject to instrument
noise, micro-scale patchiness, and tidal state variability. \code{smooth\_suitability()}
applies a Gaussian distance-weighted kernel to the suitability scores, replacing
each cell's score with a weighted mean of all cells within \code{bandwidth\_m} metres.
The result is a spatially coherent surface that better reflects broad habitat
quality rather than single-cast noise.

The smoothed score is computed as:
\deqn{\hat{s}_i = \frac{\sum_j w_{ij} \cdot s_j}{\sum_j w_{ij}}}{}
where \eqn{w_{ij} = \exp(-d_{ij}^2 / (2\sigma^2))}{} and \eqn{\sigma}{} is the
Gaussian bandwidth in metres. Excluded cells (suitability = 0, excluded flag)
are down-weighted but included so that unsuitable barriers are preserved.
\end{Description}
%
\begin{Usage}
\begin{verbatim}
smooth_suitability(
  result,
  bandwidth_m = 500,
  max_radius_m = NULL,
  smooth_excluded = FALSE,
  verbose = TRUE
)
\end{verbatim}
\end{Usage}
%
\begin{Arguments}
\begin{ldescription}
\item[\code{result}] Dataframe from \code{\LinkA{predict\_oyster()}{predict.Rul.oyster}} with \code{lat}, \code{lon},
\code{suitability}, and \code{excluded} columns.

\item[\code{bandwidth\_m}] Numeric. Gaussian kernel bandwidth (1 sigma) in metres.
Default 500 m. Increase for sparser surveys; decrease to preserve fine detail.

\item[\code{max\_radius\_m}] Numeric. Maximum search radius for neighbours in metres.
Cells further than this are ignored (default \code{3 * bandwidth\_m}).

\item[\code{smooth\_excluded}] Logical. If FALSE (default), excluded cells do not
contribute to neighbours' smoothed scores (hard barriers are preserved).

\item[\code{verbose}] Logical. Print summary (default TRUE).
\end{ldescription}
\end{Arguments}
%
\begin{Value}
Input dataframe with \code{suitability\_raw} (original) and
\code{suitability} (smoothed, overwrites original) columns, plus a
\code{suitability\_class} column reclassified from the smoothed score.
\end{Value}
%
\begin{Examples}
\begin{ExampleCode}
## Not run: 
result <- predict_oyster(survey, "ostrea_edulis")

# Smooth with 300 m bandwidth (good for dense ADCP surveys)
result_smooth <- smooth_suitability(result, bandwidth_m = 300)

# Compare raw vs smoothed
plot(result$suitability_raw, result$suitability,
     xlab = "Raw", ylab = "Smoothed", pch = 20)

## End(Not run)
\end{ExampleCode}
\end{Examples}
\HeaderA{spatial\_block\_cv}{Spatial block cross-validation for suitability model evaluation}{spatial.Rul.block.Rul.cv}
%
\begin{Description}
Evaluates model performance using spatial block cross-validation (Roberts
et al. 2017), which avoids inflated performance estimates that arise when
nearby train and test records share similar environments.

Records are partitioned into \code{n\_blocks} contiguous geographic blocks. In
each fold, one block is held out as test data and the remaining blocks
provide training data to re-score the suitability model. AUC, TSS, and
Brier score are computed per fold and summarised across folds.

This function is purely evaluative \bsl{}u2014 it does not refit a new model. Instead,
it uses the spatial heterogeneity in the existing \code{predicted} suitability
surface to estimate out-of-block performance.
\end{Description}
%
\begin{Usage}
\begin{verbatim}
spatial_block_cv(
  predicted,
  records,
  presence_col = "presence",
  n_blocks = 5L,
  match_radius_deg = 0.002,
  seed = 42L,
  plot = TRUE,
  verbose = TRUE
)
\end{verbatim}
\end{Usage}
%
\begin{Arguments}
\begin{ldescription}
\item[\code{predicted}] Dataframe from \code{\LinkA{predict\_oyster()}{predict.Rul.oyster}} containing \code{lat}, \code{lon},
\code{suitability}.

\item[\code{records}] Dataframe of known presence/absence records. Must contain
\code{lat}, \code{lon}, and a presence/absence column.

\item[\code{presence\_col}] Character. Name of presence/absence column (default
\code{"presence"}).

\item[\code{n\_blocks}] Integer. Number of spatial blocks (default 5). More blocks
= finer spatial resolution but smaller test sets per fold.

\item[\code{match\_radius\_deg}] Numeric. Matching radius in degrees (default 0.002
\bsl{}u2248 220 m). Passed to inner matching step.

\item[\code{seed}] Integer. Random seed for reproducibility (default 42).

\item[\code{verbose}] Logical. Print fold-level diagnostics (default TRUE).
\end{ldescription}
\end{Arguments}
%
\begin{Value}
A named list:
\begin{itemize}

\item{} \code{auc\_mean}: mean AUC across blocks (primary reporting metric)
\item{} \code{auc\_sd}: standard deviation of per-block AUC
\item{} \code{tss\_mean}, \code{tss\_sd}: mean / SD of True Skill Statistic
\item{} \code{brier\_mean}, \code{brier\_sd}: mean / SD of Brier score
\item{} \code{n\_blocks\_actual}: number of blocks with \bsl{}u2265 5 records (used in CV)
\item{} \code{block\_results}: dataframe with per-fold metrics
\item{} \code{spatial\_autocorr\_range\_deg}: estimated autocorrelation range (degrees)
\item{} \code{spatial\_bias\_warning}: logical \bsl{}u2014 TRUE if random CV would be misleading

\end{itemize}

\end{Value}
%
\begin{Section}{Interpreting results}

Compare \code{auc\_mean} from spatial CV with standard AUC from
\code{\LinkA{validate\_against\_records()}{validate.Rul.against.Rul.records}}. A large drop (> 0.10) indicates strong
spatial autocorrelation and means the model is less transferable than
the standard validation suggested. This is common for coastal surveys
with spatially clustered records.
\end{Section}
%
\begin{References}
Roberts et al. (2017) Ecography 40:913-929.
doi:10.1111/ecog.02881
\end{References}
%
\begin{Examples}
\begin{ExampleCode}
## Not run: 
records <- read.csv("nbn_ostrea_edulis.csv")
result  <- predict_oyster(survey, "ostrea_edulis")

# Standard validation
std_val <- validate_against_records(result, records)
std_val$auc  # may be optimistically high

# Spatial block CV
sp_cv <- spatial_block_cv(result, records, n_blocks = 5)
sp_cv$auc_mean  # more honest estimate of transferability

# Compare \u2014 if sp_cv$auc_mean << std_val$auc, model is spatially
# autocorrelated and transfers less well than initially apparent.

## End(Not run)
\end{ExampleCode}
\end{Examples}
\HeaderA{stack\_surveys}{Stack multiple survey runs from the same sensor type}{stack.Rul.surveys}
%
\begin{Description}
Combines two or more dataframes from the same sensor (e.g. an ADCP run on
Monday and again on Thursday) by row-binding them and re-averaging
overlapping cells onto a common spatial grid. Non-overlapping cells are
kept as-is.

This is the right tool when you have repeated surveys of the same area \bsl{}u2014
rather than picking one run, overlapping cells are averaged across all runs
which improves noise reduction. Non-overlapping cells (different survey
extents) are preserved.

After stacking, pass the result into \code{\LinkA{merge\_sensor\_data()}{merge.Rul.sensor.Rul.data}} alongside your
other sensor datasets as usual.
\end{Description}
%
\begin{Usage}
\begin{verbatim}
stack_surveys(..., spatial_res = 4L, verbose = TRUE)
\end{verbatim}
\end{Usage}
%
\begin{Arguments}
\begin{ldescription}
\item[\code{...}] Two or more dataframes of the same sensor type. Each must contain
\code{lat} and \code{lon} columns.

\item[\code{spatial\_res}] Integer. Decimal places for the common spatial grid
(default \code{4}, \bsl{}u224811 m at 56\bsl{}u00b0N).

\item[\code{verbose}] Logical. Print a summary (default \code{TRUE}).
\end{ldescription}
\end{Arguments}
%
\begin{Value}
A single spatially averaged dataframe combining all input surveys.
\end{Value}
%
\begin{Examples}
\begin{ExampleCode}
## Not run: 
# Stack two ADCP runs before merging with other sensors
adcp_mon <- read_nortek_adcp("survey_monday.csv")
adcp_thu <- read_nortek_adcp("survey_thursday.csv")
adcp_all <- stack_surveys(adcp_mon, adcp_thu)

survey <- merge_sensor_data(adcp = adcp_all, bathy = bathy_df)

## End(Not run)
\end{ExampleCode}
\end{Examples}
\HeaderA{tidal\_port\_info}{Look up approximate tidal range for major European survey ports}{tidal.Rul.port.Rul.info}
%
\begin{Description}
Returns the mean spring tidal range (MHWS - MLWS) and typical Chart Datum
offset from Mean Sea Level for a named port. Useful for a quick sanity check
on tidal height values before applying \code{\LinkA{correct\_to\_chart\_datum()}{correct.Rul.to.Rul.chart.Rul.datum}}.

This is a reference table only \bsl{}u2014 always use official tide table predictions
for actual corrections.
\end{Description}
%
\begin{Usage}
\begin{verbatim}
tidal_port_info(port)
\end{verbatim}
\end{Usage}
%
\begin{Arguments}
\begin{ldescription}
\item[\code{port}] Character. Port name (partial, case-insensitive match).
\end{ldescription}
\end{Arguments}
%
\begin{Value}
A one-row dataframe with columns \code{port}, \code{country},
\code{mhws\_m}, \code{mlws\_m}, \code{spring\_range\_m}, \code{cd\_below\_msl\_m}.
\end{Value}
%
\begin{Examples}
\begin{ExampleCode}
tidal_port_info("oban")
tidal_port_info("brest")
\end{ExampleCode}
\end{Examples}
\HeaderA{update\_species\_tolerances}{Update species tolerance parameters from field observations (Bayesian)}{update.Rul.species.Rul.tolerances}
%
\begin{Usage}
\begin{verbatim}
update_species_tolerances(
  records,
  species,
  update_vars = NULL,
  presence_col = "presence",
  method = c("map", "mcmc"),
  n_iter = 5000L,
  n_warmup = 1000L,
  prior_sd_fraction = 0.2,
  min_records = 20L,
  verbose = TRUE
)
\end{verbatim}
\end{Usage}
%
\begin{Arguments}
\begin{ldescription}
\item[\code{records}] Dataframe of field observations. Must contain:
\begin{itemize}

\item{} \code{lat}, \code{lon} \bsl{}u2014 coordinates
\item{} \code{presence} (or name given in \code{presence\_col}) \bsl{}u2014 1 = present, 0 = absent
\item{} Columns matching the environmental variables to update (e.g. \code{temperature},
\code{salinity}, \code{depth}, etc.)

\end{itemize}


\item[\code{species}] Character. Species key (e.g. \code{"ostrea\_edulis"}). See
\code{\LinkA{list\_species()}{list.Rul.species}}.

\item[\code{update\_vars}] Character vector. Variables to update parameters for.
Defaults to all numeric scored variables present in \code{records}.

\item[\code{presence\_col}] Character. Name of presence/absence column (default
\code{"presence"}).

\item[\code{method}] Character. \code{"map"} (default, fast) or \code{"mcmc"} (full posterior).

\item[\code{n\_iter}] Integer. MCMC iterations (default 5000; ignored for MAP).

\item[\code{n\_warmup}] Integer. MCMC warm-up / burn-in (default 1000; ignored for MAP).

\item[\code{prior\_sd\_fraction}] Numeric. Prior SD as fraction of parameter range
(default 0.20 = 20\bsl{}

\bsl{}itemmin\_recordsInteger. Minimum matched records required (default 20).

\bsl{}itemverboseLogical. Print progress and diagnostics (default TRUE).

\end{ldescription}

Invisibly: a named list with elements:
\begin{itemize}

\item{} \code{species}: species key
\item{} \code{method}: "map" or "mcmc"
\item{} \code{n\_records}: number of matched records used
\item{} \code{loglik\_null}: log-likelihood of intercept-only model
\item{} \code{loglik\_fit}: log-likelihood at MAP/posterior mean
\item{} \code{mcfadden\_r2}: 1 - loglik\_fit / loglik\_null (pseudo-R\bsl{}u00b2)
\item{} \code{updated\_params}: named list of updated parameter values per variable
\item{} \code{posterior\_sd}: named list of posterior SD per parameter (from Laplace/MCMC)
\item{} \code{prior\_params}: the parameter values before updating (for comparison)
\item{} \code{convergence}: optim() convergence code (0 = success; MAP only)

\end{itemize}


Side effect: updates the in-session cache accessed by \code{\LinkA{score\_locations()}{score.Rul.locations}}.


Fits a Bayesian logistic regression linking AHP suitability scores to
observed presence/absence, then updates the optimal-range tolerance
parameters for one or more scored variables.

\strong{Default method (MAP + Laplace):} finds the Maximum A Posteriori estimate
via \code{optim()} (L-BFGS-B) and approximates the posterior as a multivariate
normal using the numerical Hessian. Fast; suitable for routine use.

\strong{MCMC method:} Random-Walk Metropolis-Hastings sampler for full posterior
samples. Recommended when n < 50 or when credible intervals are needed for
publication. May take several minutes for large parameter sets.

Posteriors are stored in a package-level cache so that repeated calls
accumulate evidence (sequential Bayesian updating). Use
\code{\LinkA{save\_tolerance\_update()}{save.Rul.tolerance.Rul.update}} to persist between sessions.

What gets updated

For each variable in \code{update\_vars}, the function estimates shifts in:
\begin{itemize}

\item{} \code{optimal\_min}, \code{optimal\_max} \bsl{}u2014 the sweet-spot bounds
\item{} \code{acceptable\_min} / \code{poor\_min}, \code{acceptable\_max} / \code{poor\_max} \bsl{}u2014 shoulder
bounds (updated proportionally to maintain curve shape)

\end{itemize}


Parameters are constrained to biologically plausible ranges (the prior
uncertainty bounds) via box constraints in the optimiser.



## Not run: 
\# Field records with presence/absence + environmental measurements
records <- data.frame(
  lat      = c(51.5, 51.6, 51.7, 51.8, 51.4),
  lon      = c(-4.1, -4.2, -4.0, -4.3, -4.0),
  presence = c(1, 1, 0, 0, 1),
  temperature = c(16.2, 17.1, 12.3, 11.0, 15.8),
  salinity    = c(32, 33, 31, 30, 34),
  depth       = c(5, 8, 15, 20, 3)
)

fit <- update\_species\_tolerances(
  records = records,
  species = "ostrea\_edulis",
  update\_vars = c("temperature", "salinity", "depth")
)

\# See updated parameters
get\_tolerance\_posteriors("ostrea\_edulis")

\# Re-run predict\_oyster \bsl{}u2014 it will automatically use updated parameters
result <- predict\_oyster(survey, "ostrea\_edulis")

\# Save to disk for next session
save\_tolerance\_update("ostrea\_edulis")

## End(Not run)


\code{\LinkA{get\_tolerance\_posteriors()}{get.Rul.tolerance.Rul.posteriors}}, \code{\LinkA{reset\_tolerance\_update()}{reset.Rul.tolerance.Rul.update}},
\code{\LinkA{save\_tolerance\_update()}{save.Rul.tolerance.Rul.update}}, \code{\LinkA{load\_tolerance\_update()}{load.Rul.tolerance.Rul.update}}

\end{Arguments}
\HeaderA{validate\_against\_records}{Validate suitability predictions against known presence/absence records}{validate.Rul.against.Rul.records}
%
\begin{Description}
Compares the suitability scores from \code{\LinkA{predict\_oyster()}{predict.Rul.oyster}} against a
presence/absence dataset to quantify how well the model discriminates
between occupied and unoccupied habitat. Useful for assessing whether the
AHP-weighted scoring produces ecologically defensible outputs before using
them for site selection.

\strong{Metrics computed:}
\begin{itemize}

\item{} \strong{AUC} (Area Under the ROC Curve): probability that a random presence
location scores higher than a random absence location. 0.5 = random;
1.0 = perfect discrimination.
\item{} \strong{Optimal threshold:} Youden's J statistic (maximises sensitivity +
specificity - 1). Used for the confusion-matrix-derived metrics below.
\item{} \strong{Sensitivity} (true positive rate): fraction of presences correctly
predicted above threshold.
\item{} \strong{Specificity} (true negative rate): fraction of absences correctly
predicted below threshold.
\item{} \strong{F1 score}: harmonic mean of precision and sensitivity.
\item{} \strong{Brier score}: mean squared error between predicted probability and
observed presence/absence; 0 = perfect, 0.25 = uninformative.
\item{} \strong{TSS} (True Skill Statistic = sensitivity + specificity - 1): values
above 0.4 generally indicate useful predictive skill.

\end{itemize}

\end{Description}
%
\begin{Usage}
\begin{verbatim}
validate_against_records(
  predicted,
  records,
  presence_col = "presence",
  match_radius_deg = 0.002,
  plot = TRUE,
  verbose = TRUE
)
\end{verbatim}
\end{Usage}
%
\begin{Arguments}
\begin{ldescription}
\item[\code{predicted}] Dataframe from \code{\LinkA{predict\_oyster()}{predict.Rul.oyster}} containing \code{lat}, \code{lon},
and \code{suitability} columns.

\item[\code{records}] Dataframe of known presence/absence records. Must contain:
\begin{itemize}

\item{} Coordinate columns (\code{lat}/\code{lon} or equivalents)
\item{} A presence column: 1 = present, 0 = absent (or logical TRUE/FALSE).

\end{itemize}


\item[\code{presence\_col}] Character. Name of the presence/absence column in
\code{records} (default \code{"presence"}).

\item[\code{match\_radius\_deg}] Numeric. Radius in decimal degrees within which a
prediction cell is matched to a record (default \code{0.002} \bsl{}u2248 220 m).
Records that fall outside all prediction cells at this radius are dropped
with a warning.

\item[\code{plot}] Logical. Print an ASCII ROC curve to the console (default \code{TRUE}).

\item[\code{verbose}] Logical. Print full metric summary (default \code{TRUE}).
\end{ldescription}
\end{Arguments}
%
\begin{Value}
A named list:
\begin{itemize}

\item{} \code{auc}: numeric \LinkA{0,1}{0,1}
\item{} \code{optimal\_threshold}: numeric \LinkA{0,1}{0,1} \bsl{}u2014 Youden's J
\item{} \code{sensitivity}: numeric \LinkA{0,1}{0,1} at optimal threshold
\item{} \code{specificity}: numeric \LinkA{0,1}{0,1} at optimal threshold
\item{} \code{f1}: numeric \LinkA{0,1}{0,1} at optimal threshold
\item{} \code{tss}: numeric \LinkA{-1,1}{.Rdash.1,1} at optimal threshold
\item{} \code{brier\_score}: numeric \LinkA{0,1}{0,1}
\item{} \code{n\_presences}: integer
\item{} \code{n\_absences}: integer
\item{} \code{n\_unmatched}: integer \bsl{}u2014 records dropped due to no nearby prediction
\item{} \code{roc\_df}: dataframe with columns \code{threshold}, \code{sensitivity},
\code{specificity} (for custom plotting)

\end{itemize}

\end{Value}
%
\begin{Examples}
\begin{ExampleCode}
## Not run: 
# Load known flat oyster records (e.g. from NBN Atlas CSV export)
records <- read.csv("nbn_ostrea_edulis.csv")

# Run prediction on same area
result <- predict_oyster(survey, "ostrea_edulis")

# Validate
val <- validate_against_records(result, records,
                                 presence_col = "presence")
val$auc           # overall discrimination power
val$tss           # skill score
val$roc_df        # data for custom ggplot

## End(Not run)
\end{ExampleCode}
\end{Examples}
\printindex{}
\end{document}
